\documentclass[a4paper,11pt]{article}
\usepackage[T1]{fontenc}
\usepackage[utf8]{inputenc}
\usepackage[ngerman]{babel}
\usepackage{amsmath}
\usepackage{amssymb}
\usepackage[margin=2.5cm]{geometry}
\usepackage{enumitem}
\usepackage{booktabs}
\usepackage{array}
\newcolumntype{L}[1]{>{\raggedright\arraybackslash}p{#1}}
\usepackage{tikz}
\usetikzlibrary{arrows.meta,positioning,babel}
\usepackage{listings}
\usepackage{xcolor}
\usepackage[breakable]{tcolorbox}
\usepackage{hyperref}
\hypersetup{colorlinks=true, linkcolor=blue!50!black, urlcolor=blue!50!black,
  pdftitle={Webtechnologien - Erklaerung Uebungsblatt 2}}

\setlength{\parindent}{0pt}
\setlength{\parskip}{0.5em}
\setlength{\emergencystretch}{3em}

% ---------- Boxen ----------
\newtcolorbox{merkbox}[1][]{colback=blue!5, colframe=blue!40, fonttitle=\bfseries, title={Merke}, breakable, #1}
\newtcolorbox{analogie}[1][]{colback=green!5, colframe=green!40, fonttitle=\bfseries, title={Analogie}, breakable, #1}
\newtcolorbox{ergebnis}[1][]{colback=orange!8, colframe=orange!60!black, fonttitle=\bfseries, title={Ergebnis}, breakable, #1}
\newtcolorbox{tutor}[1][]{colback=gray!6, colframe=gray!50, fonttitle=\bfseries, title={Was könnte der Tutor fragen?}, breakable, #1}

% ---------- Listings ----------
\lstdefinelanguage{JavaScript}{
  keywords={const,let,var,function,return,if,else,while,for,of,try,catch,throw,new,typeof,true,false,null,undefined,this},
  sensitive=true,
  comment=[l]{//},
  morecomment=[s]{/*}{*/},
  morecomment=[s]{<!--}{-->},
  morestring=[b]",
  morestring=[b]',
  morestring=[b]`
}
\lstset{
  language=JavaScript,
  basicstyle=\ttfamily,
  keywordstyle=\color{blue!70!black}\bfseries,
  commentstyle=\color{green!45!black}\itshape,
  stringstyle=\color{red!60!black},
  showstringspaces=false,
  columns=fullflexible,
  keepspaces=true,
  upquote=true,
  breaklines=true,
  breakatwhitespace=false,
  postbreak=\mbox{\textcolor{gray}{$\hookrightarrow$}\space},
  literate={ä}{{\"a}}1 {ö}{{\"o}}1 {ü}{{\"u}}1 {Ä}{{\"A}}1 {Ö}{{\"O}}1 {Ü}{{\"U}}1 {ß}{{\ss}}1
}
\lstdefinestyle{block}{
  basicstyle=\ttfamily\small,
  frame=single, rulecolor=\color{gray!60},
  backgroundcolor=\color{gray!4},
  numbers=left, numberstyle=\tiny\color{gray}, numbersep=6pt,
  xleftmargin=1.8em, framexleftmargin=1.5em
}
\lstdefinestyle{out}{
  language={}, basicstyle=\ttfamily\small,
  frame=single, rulecolor=\color{orange!50},
  backgroundcolor=\color{orange!4},
  numbers=none, xleftmargin=0.5em, framexleftmargin=0.3em
}
\lstnewenvironment{code}[1][]{\lstset{style=block,#1}}{}
\lstnewenvironment{htmlcode}[1][]{\lstset{style=block,language=HTML,#1}}{}
\lstnewenvironment{ausgabe}[1][]{\lstset{style=out,#1}}{}

% ---------- Hilfsbefehle ----------
\newcommand{\nl}{\ensuremath{\hookleftarrow}}
\newcommand{\ws}{\textvisiblespace}
\newcommand{\qq}[1]{\textquotedbl#1\textquotedbl}

\title{\textbf{Webtechnologien -- Übungsblatt 2}\\[0.3em]
\large JavaScript \& DOM -- ausführliche Lern-Erklärung zu allen Aufgaben}
\author{}
\date{}

\begin{document}
\maketitle
\vspace{-2em}
\tableofcontents
\newpage

% =====================================================================
\section{Einstieg: Worum geht es auf diesem Blatt?}
% =====================================================================

Auf Blatt~1 ging es um HTML und CSS, also darum, \emph{wie} eine Seite aufgebaut ist und \emph{wie} sie aussieht.
Auf Blatt~2 kommt jetzt \textbf{JavaScript} (kurz JS) dazu: die Programmiersprache, mit der eine Webseite
auf den Benutzer \emph{reagieren} kann -- zum Beispiel ein Fenster öffnen, einen Preis ausrechnen oder
ein Bild austauschen, sobald jemand auf einen Knopf klickt. Dazu kommt das \textbf{DOM}, das ist die Art, wie der
Browser die HTML-Seite intern als Baum speichert, damit JavaScript sie lesen und verändern kann.
Beides wird in Teil~\ref{sec:grundlagen} von Grund auf erklärt.

Wichtig aus der Angabe: Es darf \textbf{nur reines JavaScript und CSS} verwendet werden, also keine fertigen
Bibliotheken oder Frameworks (wie jQuery oder React).

\subsection*{Wie viel muss ich machen?}

Man muss \textbf{60\,\% der Aufgaben} machen. Bei 6 Aufgaben sind das $0{,}6 \cdot 6 = 3{,}6$, aufgerundet also
\textbf{4 von 6 Aufgaben}. Bei der Kreuzerl-Übung kreuzt du nur an, was du auch wirklich erklären kannst.

\begin{center}
\begin{tabular}{@{}clll@{}}
\toprule
Rang (schwer $\to$ leicht) & Aufgabe & Thema & Punkte\\
\midrule
1 & 4 & Glaspreis (Formular + Rechnung) & 0,2\\
2 & 6 & Bildergalerie (Vor/Zurück-Buttons) & 0,2\\
3 & 3 & DOM-Operationen (Baum + 16 Ausgaben) & 0,2\\
4 & 2 & Zahlenratespiel & 0,1\\
5 & 5 & Event Listener und Target & 0,2\\
6 & 1 & Moodle-Quiz & 0,1\\
\bottomrule
\end{tabular}
\end{center}

\begin{merkbox}[title={Strategie}]
\textbf{Empfohlener Plan: Aufgaben 1, 2, 3 und 5 ankreuzen.} Das sind die vier leichtesten.
Die Aufgaben 4 und 6 sind trotzdem fertig gelöst und hier erklärt -- wenn du sie verstehst, kannst du sie
zusätzlich ankreuzen, musst aber nicht.

\textbf{Aufgabe 1} (Moodle-Quiz "`JavaScript Concepts"') kann niemand für dich lösen: Du machst das Quiz selbst in
Moodle und gibst einen \textbf{Screenshot deiner Punkte} ab (Hinweise dazu stehen in \texttt{aufgabe1/LIESMICH.txt}).
Tipp: Lies vorher Teil~\ref{sec:grundlagen}, dann fällt dir das Quiz viel leichter.
\end{merkbox}

\subsection*{Wie probiere ich den Code selbst aus?}

Jede Aufgabe liegt in einem eigenen Ordner (\texttt{aufgabe2} bis \texttt{aufgabe6}). Du öffnest die
HTML-Datei per Doppelklick im Browser (z.\,B. Chrome). Mit der Taste \textbf{F12} öffnen sich die
\textbf{Entwicklerwerkzeuge} (DevTools). Dort gibt es den Reiter \textbf{Console} (Konsole): Das ist ein
Textfenster, in das JavaScript mit dem Befehl \lstinline|console.log(...)| Dinge hineinschreiben kann.
Diese Konsole brauchst du vor allem für die Aufgaben~3 und~5.

\subsection*{Aufbau dieser Erklärung}
\begin{enumerate}[nosep]
  \item \textbf{Grundlagen} (Abschnitt~\ref{sec:grundlagen}): alles, was du für die Aufgaben wissen musst, von null an.
  \item \textbf{Die Aufgaben 1--6} (Abschnitte~\ref{sec:a1}--\ref{sec:a6}): jeweils Aufgabe in einfachen Worten,
        Lösungsidee, Code Schritt für Schritt, Ergebnis und typische Tutor-Fragen.
  \item \textbf{Zusammenfassung} (Abschnitt~\ref{sec:zusammenfassung}): Ergebnistabelle und goldene Regeln.
\end{enumerate}

% =====================================================================
\section{Grundlagen}\label{sec:grundlagen}
% =====================================================================

\subsection{Was ist JavaScript und wie kommt es in die Webseite?}

Eine Webseite besteht aus drei Schichten:
\begin{itemize}[nosep]
  \item \textbf{HTML} legt fest, \emph{was} auf der Seite ist (Überschriften, Bilder, Links, Formulare).
  \item \textbf{CSS} legt fest, \emph{wie} es aussieht (Farben, Abstände, Größen).
  \item \textbf{JavaScript} legt fest, \emph{was passiert} (Reaktion auf Klicks, Rechnungen, Änderungen der Seite).
\end{itemize}

\begin{analogie}
Stell dir ein Haus vor: HTML sind die Mauern und Räume, CSS ist die Farbe an den Wänden und die Einrichtung,
JavaScript ist die Elektrik -- du drückst einen Schalter (Klick), und das Licht geht an (etwas passiert).
\end{analogie}

JavaScript schreibt man in eine eigene Datei mit der Endung \texttt{.js} und bindet sie im HTML so ein:
\begin{htmlcode}[numbers=none]
<script src="game.js" defer></script>
\end{htmlcode}

\textbf{Warum \texttt{defer}?} Der Browser liest eine HTML-Datei von oben nach unten. Ohne \texttt{defer} würde er
beim \texttt{<script>}-Tag im \texttt{<head>} \emph{sofort} anhalten und das Skript ausführen -- zu diesem Zeitpunkt
hat er den \texttt{<body>} aber noch gar nicht gelesen. Ein Skript, das z.\,B.\ einen Link im Body sucht, würde ihn
nicht finden. \texttt{defer} heißt "`aufschieben"': Das Skript wird erst ausgeführt, \textbf{nachdem die ganze Seite
gelesen wurde}. Dann existieren alle Elemente schon.

\begin{merkbox}
\lstinline|<script src="datei.js" defer></script>| im \texttt{<head>} = "`Lade das Skript, aber führe es erst aus, wenn
die ganze HTML-Seite fertig eingelesen ist."' Das verwenden alle Aufgaben dieses Blatts.
\end{merkbox}

\subsection{Variablen: \texttt{const}, \texttt{let} und \texttt{var}}

Eine \textbf{Variable} ist ein Name, unter dem sich das Programm einen Wert merkt.

\begin{analogie}
Eine Variable ist wie eine \textbf{beschriftete Schachtel}. Auf der Schachtel steht der Name (z.\,B.\ \texttt{alter}),
in der Schachtel liegt der Wert (z.\,B.\ \texttt{15}).
\begin{itemize}[nosep]
\item \texttt{let}: Schachtel mit Deckel -- du kannst den Inhalt später austauschen.
\item \texttt{const}: Schachtel, die nach dem ersten Befüllen \textbf{zugeklebt} wird -- du kannst keinen neuen
      Inhalt hineinlegen.
\item \texttt{var}: die alte Schachtel aus früheren JavaScript-Versionen -- funktioniert, hat aber ein paar
      seltsame Eigenschaften (siehe unten).
\end{itemize}
\end{analogie}

\begin{code}
const MAX_ATTEMPTS = 3;   // fixer Wert, wird nie geändert
let keepPlaying = true;   // kann sich später ändern
keepPlaying = false;      // erlaubt, weil let
MAX_ATTEMPTS = 5;         // FEHLER: TypeError, const darf man nicht neu zuweisen
\end{code}

\textbf{Gültigkeitsbereich (Scope):} Eine mit \texttt{let} oder \texttt{const} angelegte Variable gilt nur innerhalb
der geschweiften Klammern \texttt{\{ \}}, in denen sie steht (dem sogenannten \emph{Block}). Eine \texttt{var}-Variable
gilt dagegen in der ganzen Funktion (bzw.\ im ganzen Skript).

\textbf{Zwei Fehlerarten}, die unten vorkommen: Ein \textbf{SyntaxError} heißt "`der Code ist falsch
geschrieben"' -- JS merkt das schon beim Einlesen, und das \emph{ganze} Skript läuft gar nicht erst los. Ein
\textbf{TypeError} entsteht erst \emph{während} das Programm läuft, wenn eine Operation mit einem Wert nicht
erlaubt ist (z.\,B.\ eine \texttt{const} neu zuweisen oder etwas von \texttt{undefined} lesen). Das Skript bricht an
dieser Stelle ab.

\begin{merkbox}
\begin{tabular}{@{}lll@{}}
\toprule
 & neu zuweisen? & nochmal mit gleichem Namen anlegen?\\
\midrule
\texttt{const} & nein & nein (SyntaxError)\\
\texttt{let}   & ja   & nein (SyntaxError)\\
\texttt{var}   & ja   & \textbf{ja}, wird stillschweigend erlaubt\\
\bottomrule
\end{tabular}

\medskip
\textbf{Falle:} \texttt{const} klebt nur die \emph{Schachtel} zu, nicht den Inhalt. Liegt in der Schachtel ein
Objekt (ein zusammengesetzter Wert wie eine Liste oder ein HTML-Element, siehe nächste Abschnitte), darf man das
Objekt selbst verändern:
\lstinline|const liste = []; liste.push(5);| ist erlaubt, weil \texttt{liste} weiterhin dasselbe Array ist --
es wurde nur etwas hineingelegt. Genauso ist in Aufgabe~4 \lstinline|priceInput.value = "-";| erlaubt, obwohl
\texttt{priceInput} eine \texttt{const} ist.
\end{merkbox}

\subsection{Datentypen und \texttt{typeof}}

Jeder Wert in JavaScript hat einen \textbf{Typ}, also eine Art. Die wichtigsten:

\begin{center}
\begin{tabular}{@{}llp{6.5cm}@{}}
\toprule
Typ & Beispiel & Bedeutung\\
\midrule
\texttt{number}    & \lstinline|7|, \lstinline|3.5|, \lstinline|NaN| & Zahl (ganze Zahlen und Kommazahlen sind derselbe Typ)\\
\texttt{string}    & \lstinline|"Hallo"|, \lstinline|"7"| & Text, in \texttt{\textquotedbl{}...\textquotedbl{}} oder \texttt{'...'} (beides gleichwertig)\\
\texttt{boolean}   & \lstinline|true|, \lstinline|false| & Wahrheitswert: ja oder nein\\
\texttt{undefined} & \lstinline|undefined| & "`noch kein Wert"'\\
\texttt{object}    & \lstinline|[1,2]|, \lstinline|null|, HTML-Elemente & alles Zusammengesetzte\\
\bottomrule
\end{tabular}
\end{center}

Achtung: \lstinline|7| (Zahl) und \lstinline|"7"| (Text, der zufällig aus der Ziffer 7 besteht) sind \emph{verschieden}!

\textbf{Objekte, Properties und Methoden.} Ein \textbf{Objekt} ist ein Wert, der aus mehreren benannten Teilen
besteht. Diese Teile heißen \textbf{Properties} (Eigenschaften) und werden mit einem Punkt angesprochen:
\lstinline|guesses.length| ist die Property "`Länge"' der Liste \texttt{guesses}. Ist eine Property eine Funktion,
nennt man sie \textbf{Methode} und ruft sie mit Klammern auf: \lstinline|guesses.push(5)|. Jedes HTML-Element im
Browser ist so ein Objekt, mit Properties wie \lstinline|id| oder \lstinline|value| und Methoden wie
\lstinline|remove()|.

\begin{analogie}
Ein Objekt ist wie ein \textbf{Auto}: Es hat Eigenschaften (Farbe, Kilometerstand = Properties) und Dinge, die
es tun kann (hupen, bremsen = Methoden). \lstinline|auto.farbe| fragt eine Eigenschaft ab, \lstinline|auto.hupen()|
löst eine Aktion aus.
\end{analogie}

Mit \lstinline|typeof| fragt man nach dem Typ. Das Ergebnis ist selbst ein Text:
\begin{code}[numbers=none]
typeof 7          // "number"
typeof "7"        // "string"
typeof "Invalid input. ..."   // "string"
\end{code}
Genau das nutzt Aufgabe~2: Dort liefert eine Funktion entweder eine Zahl (gültiger Tipp) oder einen Fehlertext.
Mit \lstinline|typeof guess !== "number"| erkennt man, dass es der Fehlertext ist.

\subsection{Zahlen aus Text machen: \texttt{Number()}, \texttt{NaN}, \texttt{isNaN}}

Alles, was ein Benutzer eintippt, kommt als \textbf{Text} (string) im Programm an -- auch wenn er "`7"' tippt.
Zum Rechnen und Vergleichen braucht man aber eine Zahl. \lstinline|Number(...)| wandelt um:

\begin{center}
\begin{tabular}{@{}lll@{}}
\toprule
Ausdruck & Ergebnis & Warum?\\
\midrule
\lstinline|Number("7")|   & \lstinline|7|   & Text besteht aus einer Zahl\\
\lstinline|Number("abc")| & \lstinline|NaN| & keine Zahl erkennbar\\
\lstinline|Number("")|    & \lstinline|0|   & leerer Text zählt als 0 (Eigenheit von JS)\\
\lstinline|Number(null)|  & \lstinline|0|   & \texttt{null} (= "`nichts"') wird ebenfalls zu 0\\
\lstinline|Number("2.5")| & \lstinline|2.5| & Kommazahl mit Punkt\\
\bottomrule
\end{tabular}
\end{center}

\textbf{NaN} heißt "`Not a Number"' (keine Zahl). Es ist ein spezieller Wert vom Typ \texttt{number}, der
"`Umwandlung ist schiefgegangen"' bedeutet. Seltsam, aber wichtig: \lstinline|NaN === NaN| ist \lstinline|false|
-- NaN ist mit nichts gleich, nicht einmal mit sich selbst. Deshalb gibt es die Funktion \lstinline|isNaN(x)|, die
\lstinline|true| liefert, wenn \texttt{x} NaN ist.

\begin{merkbox}[title={Merke -- automatische Umwandlung}]
Manche Rechenzeichen wandeln Text automatisch in Zahlen um, \texttt{+} aber nicht:
\begin{itemize}[nosep]
\item \lstinline|"4" / 10| ergibt \lstinline|0.4| -- Division geht nur mit Zahlen, also wandelt JS \lstinline|"4"| in 4 um.
\item \lstinline|"4" * 2| ergibt \lstinline|8| (gleicher Grund).
\item \lstinline|"4" + 1| ergibt \lstinline|"41"| -- \textbf{Falle!} \texttt{+} mit einem Text heißt "`aneinanderhängen"'.
\end{itemize}
Darum wandelt man lieber selbst mit \lstinline|Number(...)| um, dann ist es eindeutig.
\end{merkbox}

\subsection{Texte zusammenbauen}

Mit \texttt{+} hängt man Texte aneinander (\emph{String-Verkettung}). Ist einer der beiden Teile kein Text, wird
er vorher in Text umgewandelt:
\begin{code}[numbers=none]
"Attempts made: " + 2        // "Attempts made: 2"
"Guesses: " + [3, 5]         // "Guesses: 3,5"   (Array -> Elemente mit Komma)
"target: " + buttonElement   // "target: [object HTMLButtonElement]"
\end{code}
Die letzte Zeile ist wichtig für Aufgabe~5: Ein HTML-Element, das in Text umgewandelt wird, wird nur zu
\lstinline|[object HTMLButtonElement]|, also "`ein Objekt vom Typ Button-Element"'.

Alternative Schreibweise (\emph{Template-String}) mit Backticks und \lstinline|${...}|:
\begin{code}[numbers=none]
let n = 2;
`Attempts made: ${n}`        // "Attempts made: 2"
\end{code}
Beide Varianten liefern dasselbe; der Code dieses Blatts verwendet \texttt{+}.

\subsection{Vergleiche und Bedingungen}

\begin{itemize}[nosep]
\item \lstinline|===| prüft "`gleich, und zwar auch vom gleichen Typ"'. \lstinline|5 === 5| ist \lstinline|true|,
      \lstinline|5 === "5"| ist \lstinline|false|. (\lstinline|==| würde vorher umwandeln -- das vermeidet man,
      weil es überraschende Ergebnisse liefern kann.)
\item \lstinline|!==| heißt "`nicht gleich"'.
\item \lstinline|&&| heißt "`und"' (beides muss stimmen), \lstinline!||! heißt "`oder"' (eines reicht),
      \lstinline|!| heißt "`nicht"'.
\item \lstinline|if (bedingung) { ... } else if (...) { ... } else { ... }|: Es wird \emph{genau ein} Zweig
      ausgeführt, nämlich der erste, dessen Bedingung stimmt.
\end{itemize}

\textbf{Kurzschluss-Auswertung:} Bei \lstinline|a && b| schaut JS zuerst \texttt{a} an. Ist \texttt{a} schon falsch,
kann das Ganze nicht mehr wahr werden -- \texttt{b} wird gar nicht mehr geprüft. Bei \lstinline!a || b! ist es umgekehrt:
Ist \texttt{a} wahr, wird \texttt{b} übersprungen.

\subsection{Schleifen: \texttt{while} und \texttt{for}}

Eine \textbf{Schleife} wiederholt Code, solange eine Bedingung stimmt.

\begin{analogie}
\texttt{while} ist wie "`Solange du Hunger hast, iss noch einen Bissen."' Vor jedem Bissen wird gefragt.
\texttt{for} ist wie "`Mach genau 3 Liegestütze"': Man zählt mit (1, 2, 3) und hört nach der letzten auf.
\end{analogie}

\begin{code}
let keepPlaying = true;
while (keepPlaying) {                 // vor jeder Runde: noch true?
    // ... eine Spielrunde ...
    keepPlaying = confirm("Again?");  // hier entscheidet sich, ob es weitergeht
}
\end{code}

Die \texttt{for}-Schleife hat drei Teile in der Klammer, getrennt durch Strichpunkte:
\lstinline|for (start; bedingung; schritt)|.

\begin{code}
for (let attempt = 1; attempt <= 3; attempt++) {
    console.log(attempt);
}
\end{code}

Schritt für Schritt (\lstinline|attempt++| heißt "`erhöhe um 1"'):
\begin{center}
\begin{tabular}{@{}llll@{}}
\toprule
Durchlauf & \texttt{attempt} vorher & Bedingung \texttt{attempt <= 3} & Ausgabe\\
\midrule
1 & 1 & wahr & 1\\
2 & 2 & wahr & 2\\
3 & 3 & wahr & 3\\
-- & 4 & \textbf{falsch} $\to$ Schleife endet & --\\
\bottomrule
\end{tabular}
\end{center}

Für Listen gibt es noch \lstinline|for (const x of liste) { ... }|: "`Für jedes Element \texttt{x} in der Liste
mach \dots"'. Das braucht man in Aufgabe~4 für die Radio-Buttons.

\subsection{Funktionen und \texttt{return}}

Eine \textbf{Funktion} ist ein Stück Code mit Namen, das man beliebig oft aufrufen kann.

\begin{analogie}
Eine Funktion ist wie eine \textbf{Kaffeemaschine}: Du gibst etwas hinein (Bohnen, Wasser = \emph{Parameter}),
drückst den Knopf (\emph{Aufruf}), und es kommt etwas heraus (Kaffee = \emph{Rückgabewert} mit \texttt{return}).
\end{analogie}

\begin{code}
function verdopple(zahl) {    // zahl ist der Parameter
    return zahl * 2;          // Ergebnis zurückgeben
}
let x = verdopple(21);        // Aufruf mit Argument 21 -> x ist 42
\end{code}

\begin{merkbox}
\begin{itemize}[nosep]
\item \lstinline|return| \textbf{beendet die Funktion sofort} -- Code danach in der Funktion wird nicht mehr ausgeführt.
      In Aufgabe~2 beendet \lstinline|return true;| deshalb auch die Schleife innerhalb der Funktion.
\item Funktionen, die mit \lstinline|function name() {...}| angelegt werden, kann man schon \emph{vor} der Zeile,
      in der sie stehen, aufrufen. JS liest beim Start zuerst alle solchen Funktionsdefinitionen ein
      (das nennt man \emph{Hoisting}, "`Hochziehen"'). Deshalb darf in Aufgabe~2 \texttt{playNumberGame} oben
      aufgerufen werden, obwohl die Funktion erst weiter unten steht.
\end{itemize}
\end{merkbox}

\subsection{Arrays und \texttt{push}}

Ein \textbf{Array} ist eine geordnete Liste von Werten.

\begin{analogie}
Ein Array ist wie eine \textbf{nummerierte Einkaufsliste} -- nur dass die Nummerierung bei \textbf{0} beginnt.
\lstinline|push| schreibt einen neuen Punkt unten an die Liste.
\end{analogie}

\begin{code}
let guesses = [];      // leere Liste
guesses.push(3);       // [3]
guesses.push(5);       // [3, 5]
guesses[0];            // 3   (erstes Element hat Index 0!)
guesses.length;        // 2   (Anzahl der Elemente)
\end{code}

\subsection{Dialogfenster: \texttt{alert}, \texttt{prompt}, \texttt{confirm}}

Der Browser hat drei eingebaute Mini-Fenster. Solange so ein Fenster offen ist, \textbf{wartet} das Programm.

\begin{center}
\begin{tabular}{@{}lp{3.6cm}p{5.6cm}@{}}
\toprule
Funktion & zeigt & gibt zurück\\
\midrule
\lstinline|alert("Hi")|   & Nachricht + OK-Knopf          & nichts (\texttt{undefined})\\
\lstinline|prompt("Zahl?")| & Nachricht + Eingabefeld     & eingegebener \textbf{Text}; bei "`Abbrechen"' \lstinline|null|\\
\lstinline|confirm("Nochmal?")| & Nachricht + OK/Abbrechen & \lstinline|true| (OK) oder \lstinline|false| (Abbrechen)\\
\bottomrule
\end{tabular}
\end{center}

\subsection{Fehlerbehandlung: \texttt{try}, \texttt{catch}, \texttt{throw}}

Manchmal geht etwas schief, z.\,B.\ der Benutzer tippt "`abc"' statt einer Zahl. Dafür gibt es drei Bausteine:
\begin{itemize}[nosep]
\item \lstinline|throw new Error("Text")| -- "`wirf"' einen Fehler: Der normale Ablauf bricht \emph{sofort} ab.
\item \lstinline|try { ... }| -- "`versuche"' diesen Code.
\item \lstinline|catch (error) { ... }| -- wird ein Fehler geworfen, landet das Programm hier.
      In \lstinline|error.message| steht der Text des Fehlers.
\end{itemize}

\begin{analogie}
Ein Trapezkünstler (\texttt{try}) turnt über einem \textbf{Sicherheitsnetz} (\texttt{catch}). Wenn er fällt
(\texttt{throw}), stürzt nicht die ganze Vorstellung ab, sondern das Netz fängt ihn auf, und es geht kontrolliert weiter.
\end{analogie}

\begin{code}
try {
    let x = Number("abc");                // NaN
    if (isNaN(x)) {
        throw new Error("Das ist keine Zahl");
    }
    console.log("wird übersprungen");     // kommt nicht mehr dran
} catch (error) {
    console.log(error.message);           // "Das ist keine Zahl"
}
\end{code}

\subsection{Zufallszahlen: \texttt{Math.random} und \texttt{Math.floor}}\label{sec:zufall}

\begin{itemize}[nosep]
\item \lstinline|Math.random()| liefert eine zufällige Kommazahl \textbf{von 0 (inklusive) bis 1 (exklusive)},
      also z.\,B.\ \texttt{0.0}, \texttt{0.73}, \texttt{0.9999}, aber nie genau \texttt{1}. Mathematisch: $[0, 1)$.
\item \lstinline|Math.floor(x)| rundet \textbf{immer ab} auf die nächste ganze Zahl: \texttt{Math.floor(7.9)} ist \texttt{7}.
      ("`floor"' = Fußboden: man fällt nach unten.)
\end{itemize}

\textbf{Vollständiges Beispiel} für \lstinline|Math.floor(Math.random() * 10) + 1| mit dem Zufallswert $0{,}73$:
\begin{align*}
\texttt{Math.random()} &= 0{,}73 && \text{angenommener Zufallswert}\\
0{,}73 \cdot 10 &= 7{,}3 && \text{Bereich von } [0,1) \text{ auf } [0,10) \text{ strecken}\\
\texttt{Math.floor}(7{,}3) &= 7 && \text{abrunden: nur noch ganze Zahlen } 0 \dots 9\\
7 + 1 &= 8 && \text{verschieben: aus } 0\dots 9 \text{ wird } 1 \dots 10
\end{align*}

Die Randfälle zeigen, dass wirklich genau $1$ bis $10$ herauskommt:
\begin{align*}
\text{kleinster Wert: } & 0 \cdot 10 = 0, && \lfloor 0 \rfloor + 1 = 1\\
\text{größter Wert: }   & 0{,}9999 \cdot 10 = 9{,}999, && \lfloor 9{,}999 \rfloor + 1 = 9 + 1 = 10
\end{align*}
Die 11 kann nie vorkommen, weil \lstinline|Math.random()| nie genau 1 liefert.

\begin{merkbox}
Zufällige ganze Zahl von $1$ bis $n$: \lstinline|Math.floor(Math.random() * n) + 1|.
\emph{Mal $n$} legt fest, wie viele verschiedene Werte es gibt; \emph{floor} macht ganze Zahlen daraus;
\emph{$+1$} verschiebt den Start von 0 auf 1.
\end{merkbox}

\subsection{Rechnen mit Rest: Modulo (\texttt{\%})}\label{sec:modulo}

Der \textbf{Modulo}-Operator \texttt{\%} liefert den \textbf{Rest einer Division}. Beispiel: $7 \,\%\, 3 = 1$, denn
$7 = 2 \cdot 3 + 1$ -- die 3 passt zweimal in die 7, übrig bleibt 1.

\begin{analogie}
Eine \textbf{Uhr} rechnet mit Modulo 12: Es ist 11 Uhr, in 3 Stunden ist es nicht 14 Uhr, sondern 2 Uhr --
nach der 12 geht es wieder von vorne los. $(11 + 3) \,\%\, 12 = 14 \,\%\, 12 = 2$.
Genau so ein "`Im-Kreis-Laufen"' (\emph{Wrap-around}) brauchen wir in Aufgabe~6 für die Bildergalerie:
Nach dem letzten Bild kommt wieder das erste.
\end{analogie}

\textbf{Vollständiges Beispiel} mit 3 Bildern (Indizes 0, 1, 2):
\begin{align*}
0 \,\%\, 3 &= 0 && 3 \text{ passt 0-mal in } 0, \text{ Rest } 0\\
1 \,\%\, 3 &= 1 && \text{Rest } 1\\
2 \,\%\, 3 &= 2 && \text{Rest } 2\\
3 \,\%\, 3 &= 0 && 3 = 1\cdot 3 + 0 \;\Rightarrow\; \text{wieder bei } 0 \text{ (Wrap-around!)}\\
4 \,\%\, 3 &= 1 && 4 = 1\cdot 3 + 1
\end{align*}

\begin{merkbox}[title={Merke -- Falle bei negativen Zahlen}]
In JavaScript ist \lstinline|-1 % 3| gleich \lstinline|-1| (nicht 2!). Das Ergebnis behält das Vorzeichen der linken Zahl.
Ein Index $-1$ existiert in einem Array aber nicht. \textbf{Trick:} vorher die Länge $n$ addieren:
$(-1 + 3) \,\%\, 3 = 2 \,\%\, 3 = 2$. Allgemein: \lstinline|(index + n) % n| liegt immer zwischen $0$ und $n-1$,
solange \texttt{index} nicht kleiner als $-n$ ist (bei Schritten von $\pm 1$ ist das immer erfüllt).
\end{merkbox}

% ---------------------------------------------------------------------
\subsection{Das DOM: die Webseite als Baum}\label{sec:dom}

Wenn der Browser eine HTML-Datei lädt, liest er den Text und baut daraus im Speicher eine
\textbf{Baumstruktur aus Objekten}. Der Teil des Browsers, der den HTML-Text liest und daraus den Baum baut,
heißt \textbf{Parser} ("`Zerleger"'). Der fertige Baum heißt \textbf{DOM} (\emph{Document Object Model}). JavaScript arbeitet
nie mit dem HTML-Text selbst, sondern immer mit diesem Baum. Ändert JS den Baum, ändert sich sofort, was man sieht.

\begin{analogie}
Das DOM ist wie ein \textbf{Familienstammbaum}: Ganz oben steht das Dokument (die Urgroßmutter), darunter
\texttt{<html>}, darunter \texttt{<head>} und \texttt{<body>} (Geschwister), und so weiter. Jedes Element hat
\emph{genau einen} Elternteil (\texttt{parentNode}) und beliebig viele Kinder (\texttt{childNodes}), mit einem
ersten (\texttt{firstChild}) und einem letzten Kind (\texttt{lastChild}). Jeder Eintrag im Baum heißt \textbf{Knoten}
(englisch \emph{node}).
\end{analogie}

\textbf{Vollständiges Beispiel.} Aus dem HTML \lstinline[language=HTML]|<p>Hallo <b>Welt</b></p>| macht der Browser diesen Baum:

\begin{center}
\begin{tikzpicture}[
  level distance=1.3cm, sibling distance=4cm,
  el/.style={draw=blue!60, fill=blue!8, rounded corners, font=\ttfamily\small, inner sep=3pt},
  tx/.style={draw=orange!70, fill=orange!8, rounded corners, font=\ttfamily\small, inner sep=3pt},
  edge from parent/.style={draw, gray}]
\node[el] {p (Element)}
  child {node[tx] {\#text \qq{Hallo\ws}}}
  child {node[el] {b (Element)}
     child {node[tx] {\#text \qq{Welt}}}};
\end{tikzpicture}
\end{center}

Beobachtungen: (1) Der Text "`Hallo "' ist ein eigener Knoten, ein \textbf{Textknoten}. (2) Das Leerzeichen nach
"`Hallo"' gehört zum Textknoten. (3) \texttt{<b>} ist ein Kind von \texttt{<p>}, und "`Welt"' ist ein Kind von
\texttt{<b>}, weil es \emph{innerhalb} von \texttt{<b>} steht.

\subsubsection*{Knotentypen und \texttt{nodeType}}

Jeder Knoten hat eine Eigenschaft \texttt{nodeType}, eine Nummer, die sagt, welche Art Knoten es ist:

\begin{center}
\begin{tabular}{@{}clll@{}}
\toprule
\texttt{nodeType} & Name & Beispiel & \texttt{nodeName}\\
\midrule
1 & \texttt{ELEMENT\_NODE}   & \texttt{<a>}, \texttt{<p>}, \texttt{<img>} & \texttt{\textquotedbl{}A\textquotedbl{}}, \texttt{\textquotedbl{}P\textquotedbl{}} (Großbuchstaben)\\
2 & \texttt{ATTRIBUTE\_NODE} & \texttt{href=\textquotedbl{}...\textquotedbl{}}                       & \texttt{\textquotedbl{}href\textquotedbl{}}\\
3 & \texttt{TEXT\_NODE}      & Text zwischen Tags                         & \texttt{\textquotedbl{}\#text\textquotedbl{}}\\
8 & \texttt{COMMENT\_NODE}   & \texttt{<!-- ... -->}                      & \texttt{\textquotedbl{}\#comment\textquotedbl{}}\\
9 & \texttt{DOCUMENT\_NODE}  & \texttt{document}                          & \texttt{\textquotedbl{}\#document\textquotedbl{}}\\
\bottomrule
\end{tabular}
\end{center}

Weitere Eigenschaften jedes Knotens:
\begin{itemize}[nosep]
\item \texttt{nodeName}: der Name. Bei HTML-Elementen \textbf{immer in Großbuchstaben} (\texttt{\textquotedbl{}BODY\textquotedbl{}}, \texttt{\textquotedbl{}A\textquotedbl{}}),
      egal wie es im HTML geschrieben wurde.
\item \texttt{nodeValue}: der Wert. Bei Text- und Attributknoten ist das der Text bzw.\ der Attributwert.
      Bei \textbf{Elementen ist er immer \texttt{null}} -- ein Element "`hat"' keinen eigenen Wert, sondern nur Kinder.
\end{itemize}

\subsubsection*{Whitespace-Textknoten -- die häufigste Überraschung}

\textbf{Whitespace} heißt "`weißer Raum"' und meint unsichtbare Zeichen: Leerzeichen, Tabulatoren, Zeilenumbrüche.
Der Browser wirft diese Zeichen beim Bauen des Baums \textbf{nicht weg}. Jeder Zeilenumbruch und jede Einrückung
zwischen zwei Tags wird zu einem eigenen Textknoten. Beispiel:
\begin{htmlcode}
<ul>
  <li>A</li>
</ul>
\end{htmlcode}
Zwischen \texttt{<ul>} und \texttt{<li>} stehen ein Zeilenumbruch und zwei Leerzeichen. Die Kinder von
\texttt{<ul>} sind deshalb: \texttt{\#text} "`\nl\ws\ws"', dann \texttt{<li>}, dann \texttt{\#text} "`\nl"'.
\texttt{ul.firstChild} ist also \textbf{nicht} das \texttt{<li>}, sondern der unsichtbare Textknoten!
In dieser Erklärung steht \nl\ für einen Zeilenumbruch und \ws\ für ein Leerzeichen.

Wer nur Elemente will, nimmt \texttt{firstElementChild} statt \texttt{firstChild} (bzw.\ \texttt{children}
statt \texttt{childNodes}).

\subsubsection*{Attribute sind keine Kinder}

Ein \textbf{Attribut} ist eine Zusatzangabe im Start-Tag, z.\,B.\ \texttt{href} und \texttt{target} in
\lstinline[language=HTML]|<a href="https://de.wikipedia.org" target="_new">|. Attribute \emph{hängen} am Element,
sind aber \textbf{keine Kindknoten} -- sie tauchen in \texttt{childNodes} und \texttt{firstChild} nie auf.
Man erreicht sie über eigene Wege:
\begin{itemize}[nosep]
\item \lstinline|el.attributes| -- Liste aller Attributknoten, per Name (\lstinline|el.attributes["href"]|) oder per
      Nummer (\lstinline|el.attributes[0]|, Reihenfolge wie im HTML) ansprechbar; \lstinline|.length| = Anzahl.
\item \lstinline|el.getAttributeNode("href")| -- liefert den Attributknoten mit diesem Namen.
\item \lstinline|el.getAttribute("href")| -- liefert direkt den Text des Attributs.
\end{itemize}

\begin{analogie}
Ein Element ist wie eine Person im Stammbaum. Die Kinder sind ihre Kinder. Die Attribute sind wie ihre
\textbf{Ausweisdaten} (Name, Geburtsdatum): Sie gehören zur Person, sind aber keine Kinder.
\end{analogie}

\subsubsection*{Attribut oder Property? Am Beispiel \texttt{href}}

Es gibt zwei Sichtweisen auf dieselbe Information:
\begin{itemize}[nosep]
\item Das \textbf{Attribut} ist genau das, was \textbf{im HTML-Text steht}, Zeichen für Zeichen.
\item Die \textbf{Property} (Eigenschaft) ist das, was der \textbf{Browser daraus gemacht hat}. Man liest sie mit
      Punkt-Schreibweise: \lstinline|el.href|.
\end{itemize}

Bei \texttt{href} sieht man den Unterschied: Steht im HTML \texttt{href=\textquotedbl{}https://de.wikipedia.org\textquotedbl{}}, dann liefert
das Attribut genau diesen Text. Die Property \lstinline|link.href| liefert aber
\begin{center}\texttt{\textquotedbl{}https://de.wikipedia.org/\textquotedbl{}}\end{center}
-- also mit Schrägstrich am Ende. Der Browser hat die Adresse \emph{aufgelöst und vereinheitlicht} (normalisiert):
Jede vollständige Webadresse hat einen Pfad, und der leere Pfad wird als "`/"' geschrieben.
Bei einem relativen Link wie \texttt{href=\textquotedbl{}seite.html\textquotedbl{}} wäre der Unterschied noch größer: Das Attribut wäre
\texttt{\textquotedbl{}seite.html\textquotedbl{}}, die Property die volle Adresse inklusive Ordner. (Eine
\emph{absolute} Adresse ist vollständig, mit \texttt{https://} und Servername; eine \emph{relative} Adresse ist nur
ein Weg ausgehend von der aktuellen Seite, z.\,B.\ ein Dateiname im selben Ordner.)

\begin{analogie}
Das Attribut ist der \textbf{Bestellzettel} ("`1 Pizza Margherita"'), die Property ist die \textbf{gelieferte Pizza}
(fertig gebacken, in Stücke geschnitten). Inhaltlich dasselbe, aber vom Browser bearbeitet.
\end{analogie}

% ---------------------------------------------------------------------
\subsection{Inhalte lesen: \texttt{innerHTML}, \texttt{outerHTML}, \texttt{innerText}, \texttt{outerText}, \texttt{textContent}}\label{sec:inhalt}

Fünf Properties liefern den Inhalt eines Elements als Text -- jede auf eine andere Art. Einfaches Beispiel:
\lstinline[language=HTML]|<span id="s">Hallo <b>Welt</b></span>|

\begin{center}
\begin{tabular}{@{}lp{6.6cm}p{4.6cm}@{}}
\toprule
Property & Ergebnis für \texttt{s} & was wird geliefert?\\
\midrule
\texttt{innerHTML}   & \lstinline|Hallo <b>Welt</b>| & HTML-Code \emph{innerhalb} des Elements\\
\texttt{outerHTML}   & \lstinline|<span id="s">Hallo <b>Welt</b></span>| & HTML-Code \emph{inklusive} Element selbst\\
\texttt{textContent} & \lstinline|Hallo Welt| & alle Textknoten roh hintereinander\\
\texttt{innerText}   & \lstinline|Hallo Welt| & Text so, wie er \emph{angezeigt} wird\\
\texttt{outerText}   & \lstinline|Hallo Welt| & beim Lesen wie \texttt{innerText}\\
\bottomrule
\end{tabular}
\end{center}

Bei diesem einfachen Beispiel sehen die letzten drei gleich aus. Die Unterschiede zeigen sich erst bei
Whitespace und Blockelementen (siehe Aufgabe~3). Zur Erinnerung aus HTML/CSS: Ein \textbf{Blockelement}
(z.\,B.\ \texttt{<p>}, \texttt{<div>}) beginnt immer in einer neuen Zeile und danach geht es wieder in einer neuen
Zeile weiter. Ein \textbf{Inline-Element} (z.\,B.\ \texttt{<a>}, \texttt{<b>}, \texttt{<img>}, \texttt{<span>})
fließt einfach im Text mit, wie ein Wort in einem Satz. Außerdem gilt beim Anzeigen: Mehrere Leerzeichen und
Zeilenumbrüche hintereinander werden zu \emph{einem} Leerzeichen zusammengefasst.
\begin{itemize}[nosep]
\item \texttt{textContent} nimmt den Text \textbf{roh}, Zeichen für Zeichen wie im Quelltext, mit allen
      Zeilenumbrüchen und Einrückungen. Es schaut nicht darauf, wie die Seite aussieht.
\item \texttt{innerText} liefert den Text so, wie ein Mensch ihn \textbf{auf dem Bildschirm sieht}: Mehrere
      Leerzeichen und Zeilenumbrüche werden (wie beim Anzeigen) zu einem Leerzeichen zusammengefasst, und nach
      Blockelementen wie \texttt{<p>} kommen Zeilenumbrüche, weil der Absatz auf dem Bildschirm eine eigene Zeile hat.
\item \texttt{outerText} ist beim \emph{Lesen} identisch mit \texttt{innerText}. Der Unterschied ist nur beim
      \emph{Setzen}: \lstinline|el.innerText = "x"| ersetzt den Inhalt, \lstinline|el.outerText = "x"| ersetzt das
      ganze Element durch den Text.
\item Bilder haben keinen Text. Ihr \texttt{alt}-Text ist ein Attribut, kein Textknoten -- er kommt in
      \texttt{textContent} und \texttt{innerText} deshalb nicht vor.
\end{itemize}

\begin{analogie}
Stell dir ein \textbf{Geschenkpaket} vor. \texttt{outerHTML} ist das Paket \emph{mit} Verpackung und Etikett.
\texttt{innerHTML} ist der Inhalt mit allen inneren Schachteln und deren Beschriftungen. \texttt{textContent} ist
ein Zettel, auf dem jemand stur alle Texte abgeschrieben hat, inklusive Leerzeilen. \texttt{innerText} ist das, was
du beim Auspacken tatsächlich \emph{siehst}.
\end{analogie}

% ---------------------------------------------------------------------
\subsection{Elemente finden, erzeugen, einfügen, löschen}\label{sec:finden}

\subsubsection*{Finden}

\begin{center}
\begin{tabular}{@{}lp{6.2cm}@{}}
\toprule
Methode & liefert\\
\midrule
\lstinline|document.getElementById("width")| & das eine Element mit \texttt{id=\textquotedbl{}width\textquotedbl{}}\\
\lstinline|document.getElementsByTagName("a")| & Liste \emph{aller} \texttt{<a>}-Elemente\\
\lstinline|document.querySelector("h1 + p")| & das \emph{erste} Element, das zum CSS-Selektor passt\\
\lstinline|document.querySelectorAll(".gallery img")| & Liste \emph{aller} passenden Elemente\\
\bottomrule
\end{tabular}
\end{center}

Bei Listen zählt man \textbf{ab 0}: \lstinline|getElementsByTagName("a")[1]| ist der \textbf{zweite} Link.
\lstinline|.item(2)| ist genau dasselbe wie \lstinline|[2]|, nur als Methode geschrieben.

\texttt{querySelector} versteht dieselben Selektoren wie CSS:
\begin{itemize}[nosep]
\item \texttt{\#button} -- Element mit \texttt{id=\textquotedbl{}button\textquotedbl{}}; \texttt{.desc} -- Element mit \texttt{class=\textquotedbl{}desc\textquotedbl{}}.
\item \texttt{\#displayDIV .bigIMG} -- ein \texttt{.bigIMG} irgendwo \emph{innerhalb} von \texttt{\#displayDIV}.
\item \texttt{h1 + p} -- ein \texttt{<p>}, das \emph{direkt nach} einem \texttt{<h1>} kommt (gleicher Elternteil).
\item \lstinline|input[name="thickness"]| -- \texttt{<input>} mit dem Attribut \texttt{name=\textquotedbl{}thickness\textquotedbl{}}.
\item \texttt{:checked} -- nur angehakte Radio-Buttons/Checkboxen.
\end{itemize}

\subsubsection*{Erzeugen, einfügen, löschen}

\begin{analogie}
\lstinline|document.createElement("button")| ist wie einen \textbf{Lego-Stein bauen}: Er existiert jetzt, liegt aber
noch neben der Bauplatte -- man sieht ihn auf der Seite nicht. Erst \lstinline|append| bzw.\ \lstinline|appendChild|
\textbf{steckt ihn an} (hängt ihn als letztes Kind an ein Element). \lstinline|remove()| nimmt einen Stein samt
allem, was darauf steckt, wieder ab.
\end{analogie}

\begin{code}
let li = document.createElement("li");   // 1. neues Element bauen (noch unsichtbar)
li.textContent = "Neuer Punkt";          // 2. Text hineinschreiben
document.querySelector("ul").appendChild(li);  // 3. als letztes Kind einhängen -> sichtbar
document.querySelector("h1 + p").remove();     // ein Element löschen
link.removeAttribute("target");          // nur ein Attribut entfernen, Element bleibt
\end{code}

\begin{itemize}[nosep]
\item \lstinline|appendChild(x)| hängt \emph{genau einen} Knoten an.
\item \lstinline|append(x, y, ...)| kann \emph{mehrere} auf einmal anhängen (auch Texte). Beides hängt ans \textbf{Ende}.
\item Eigenschaften setzt man per Punkt: \lstinline|input.id = "price"|, \lstinline|input.readOnly = true|,
      \lstinline|button.style.position = "absolute"|.
\item Beim Label heißt die Property \lstinline|htmlFor| statt \texttt{for}, weil \texttt{for} in JavaScript schon
      ein reserviertes Wort ist (die \texttt{for}-Schleife).
\item \lstinline|el.parentElement| ist der Elternteil, wenn er ein Element ist (fast immer dasselbe wie \texttt{parentNode}).
\end{itemize}

% ---------------------------------------------------------------------
\subsection{Events: auf den Benutzer reagieren}\label{sec:events}

Ein \textbf{Event} (Ereignis) ist etwas, das passiert: ein Klick (\texttt{click}), eine Eingabe in ein Feld
(\texttt{input}), eine geänderte Auswahl (\texttt{change}). Mit \lstinline|addEventListener| sagt man einem Element:
"`Wenn dieses Ereignis bei dir passiert, führe diese Funktion aus."' Die Funktion heißt dann
\textbf{Event Listener} (englisch \emph{listener} = Zuhörer) oder \emph{Event-Handler}.

\begin{analogie}
Ein Event Listener ist wie eine \textbf{Türklingel}: Du montierst sie einmal (\texttt{addEventListener}), und jedes Mal,
wenn jemand drückt (Event), läutet es (die Funktion läuft). Du musst nicht ständig selbst nachschauen, ob jemand
vor der Tür steht.
\end{analogie}

\begin{code}
widthInput.addEventListener("input", update);  // update OHNE Klammern!
\end{code}
\textbf{Warum ohne Klammern?} \lstinline|update| ist die Funktion selbst ("`hier ist die Klingel-Melodie"').
\lstinline|update()| würde die Funktion \emph{sofort jetzt} ausführen und nur ihr Ergebnis übergeben -- das ist falsch.

\textbf{\texttt{input} oder \texttt{change}?} \texttt{input} feuert bei \emph{jeder} Änderung, also bei jedem
Tastendruck. \texttt{change} feuert, wenn die Änderung "`abgeschlossen"' ist -- bei einem Radio-Button genau dann,
wenn er ausgewählt wird.

\subsubsection*{Das Event-Objekt: \texttt{target}, \texttt{currentTarget}, \texttt{this}}

Die Listener-Funktion bekommt automatisch ein \textbf{Event-Objekt} mit Informationen übergeben (meist
\texttt{event} oder \texttt{e} genannt):
\begin{itemize}[nosep]
\item \lstinline|event.target| -- das Element, auf dem das Ereignis \textbf{ausgelöst} wurde, also das tatsächlich
      angeklickte, \emph{innerste} Element. Bleibt auf dem ganzen Weg gleich.
\item \lstinline|event.currentTarget| -- das Element, an dem der \textbf{gerade laufende Listener hängt}.
\item \lstinline|this| -- in einer normalen \lstinline|function| als Listener dasselbe wie \texttt{currentTarget}.
      (Bei einer Pfeilfunktion \lstinline|(e) => {...}| wäre \texttt{this} \emph{nicht} das Element, sondern
      käme aus dem umgebenden Code.)
\end{itemize}

\subsubsection*{Bubbling: Events steigen nach oben}

Klickt man auf ein Element, bekommen \textbf{auch alle seine Vorfahren} den Klick mit -- der Reihe nach von innen
nach außen: zuerst das Element selbst, dann sein Elternteil, dann dessen Elternteil, \dots, bis zu \texttt{document}.
Das nennt man \textbf{Bubbling} (aufsteigen wie eine Luftblase im Wasser). Jeder Listener auf diesem Weg wird ausgeführt.
(Genau genommen läuft das Event vorher auch einmal von oben nach unten -- die \emph{Capturing}-Phase --, aber
\texttt{addEventListener} hört standardmäßig nur auf die Bubbling-Phase. Für dieses Blatt reicht Bubbling.)

\begin{analogie}
Bubbling ist wie eine \textbf{Meldung in einer Firma}: Ein Mitarbeiter (Button) meldet ein Problem. Er kümmert
sich selbst darum, dann geht die Meldung an seine Abteilungsleiterin (div), dann an den Bereichsleiter (body) und
bis zur Chefin (document). Jede Person auf dem Weg kann reagieren. Auf der Meldung steht immer, \emph{wer} sie
ursprünglich geschrieben hat (\texttt{target}) -- aber wer sie \emph{gerade} liest (\texttt{currentTarget}), ändert
sich bei jeder Station.
\end{analogie}

\textbf{Vollständiges Beispiel:} \lstinline[language=HTML]|<ul id="liste"><li>Apfel</li></ul>|, ein Listener hängt
nur am \texttt{<ul>}. Man klickt auf "`Apfel"':
\begin{enumerate}[nosep]
\item Der Klick passiert auf dem \texttt{<li>}. Dort hängt kein Listener, es passiert nichts.
\item Das Event steigt zum \texttt{<ul>} auf. Dort hängt der Listener und läuft.
\item In diesem Listener gilt:\\
      \texttt{event.target} ist das \texttt{<li>}, denn dort wurde geklickt;\\
      \texttt{event.currentTarget} ist das \texttt{<ul>}, denn dort hängt der Listener.
\end{enumerate}

\subsubsection*{\texttt{stopPropagation}: Aufsteigen stoppen}

\lstinline|event.stopPropagation()| sagt: "`Ab hier nicht mehr weiterreichen."' Das Event erreicht dann
\textbf{keine weiteren Elemente} weiter oben. Weitere Listener am \emph{selben} Element laufen aber trotzdem noch
(die würde nur \lstinline|stopImmediatePropagation()| stoppen).

\begin{merkbox}
\begin{itemize}[nosep]
\item \texttt{target} = \emph{wo} geklickt wurde (bleibt gleich). \texttt{currentTarget} = \emph{wer} gerade
      reagiert (wandert mit). \texttt{this} = wie \texttt{currentTarget} (bei normaler \texttt{function}).
\item Bubbling: Event läuft vom geklickten Element nach oben bis \texttt{document}; alle Listener unterwegs laufen.
\item \texttt{stopPropagation()} stoppt den Weg nach oben, nicht die anderen Listener am selben Element.
\item Nicht verwechseln: \texttt{preventDefault()} verhindert das \emph{Standardverhalten} des Browsers (z.\,B.\
      dass ein Link geöffnet wird) -- das ist etwas anderes als das Aufsteigen zu stoppen.
\end{itemize}
\end{merkbox}

% ---------------------------------------------------------------------
\subsection{Formulare und Validierung}\label{sec:formulare}

HTML-Formularfelder können selbst Regeln haben, die der Browser prüft:
\begin{itemize}[nosep]
\item \texttt{required} -- Feld muss ausgefüllt sein.
\item \texttt{min=\textquotedbl{}0\textquotedbl{}}, \texttt{max=\textquotedbl{}200\textquotedbl{}} -- kleinster und größter erlaubter Wert (bei \texttt{type=\textquotedbl{}number\textquotedbl{}}).
\item \texttt{step=\textquotedbl{}1\textquotedbl{}} -- nur Werte in Schritten von 1, ausgehend von \texttt{min}, also nur ganze Zahlen.
\end{itemize}

\lstinline|feld.checkValidity()| fragt: "`Erfüllt dieses Feld gerade alle seine Regeln?"' und liefert
\lstinline|true| oder \lstinline|false|. Der Browser markiert ungültige Felder außerdem mit der CSS-Pseudoklasse
\texttt{:invalid} (in Aufgabe~4 werden sie dadurch rot).

\textbf{Vollständiges Beispiel} für \lstinline[language=HTML]|<input type="number" min="0" max="200" step="1" required>|:
\begin{center}
\begin{tabular}{@{}lll@{}}
\toprule
Eingabe & \texttt{checkValidity()} & Grund\\
\midrule
(leer)  & \lstinline|false| & \texttt{required} verletzt\\
\texttt{150} & \lstinline|true|  & alle Regeln erfüllt\\
\texttt{300} & \lstinline|false| & größer als \texttt{max=200}\\
\texttt{2.5} & \lstinline|false| & kein ganzzahliger Schritt (\texttt{step=1})\\
\texttt{-1}  & \lstinline|false| & kleiner als \texttt{min=0}\\
\bottomrule
\end{tabular}
\end{center}

\textbf{Radio-Gruppen:} Radio-Buttons mit demselben \texttt{name} bilden eine Gruppe; man kann nur einen davon
auswählen. Steht \texttt{required} an einem Radio der Gruppe, gilt das für die \emph{ganze Gruppe}. Dann ist
\textbf{jedes} Radio dieser Gruppe genau dann gültig, wenn \textbf{irgendein} Radio der Gruppe ausgewählt ist.

\begin{merkbox}
Der Wert eines Eingabefelds (\lstinline|feld.value|) ist \textbf{immer ein Text}, auch bei \texttt{type=\textquotedbl{}number\textquotedbl{}}.
Zum Rechnen umwandeln! Ein Farbfeld (\texttt{type=\textquotedbl{}color\textquotedbl{}}) hat immer einen Wert der Form \texttt{\textquotedbl{}\#rrggbb\textquotedbl{}}
in Kleinbuchstaben, z.\,B.\ \texttt{\textquotedbl{}\#ff0000\textquotedbl{}} (Standardwert: \texttt{\textquotedbl{}\#000000\textquotedbl{}}, Schwarz).
\end{merkbox}

\subsubsection*{Hex-Farben in 1 Minute}

Eine Farbe \texttt{\#rrggbb} besteht aus drei Paaren: Rot-, Grün- und Blauanteil. Jedes Paar ist eine Zahl im
\textbf{Sechzehnersystem} (hexadezimal, Ziffern \texttt{0}--\texttt{9} und \texttt{a}--\texttt{f}) von
\texttt{00} (= 0, nichts) bis \texttt{ff} (= 255, volle Stärke). Beispiel: \texttt{80} hexadezimal ist
$8 \cdot 16 + 0 = 128$, also halbe Stärke.

Ein \textbf{Grauton} entsteht, wenn alle drei Anteile \textbf{gleich} sind ($R = G = B$): \texttt{\#000000} ist Schwarz,
\texttt{\#808080} Mittelgrau, \texttt{\#ffffff} Weiß. Sobald ein Anteil abweicht, ist die Farbe "`bunt"', z.\,B.\
\texttt{\#ff0000} (nur Rot).

% ---------------------------------------------------------------------
\subsection{Typische Fallen auf einen Blick}

\begin{merkbox}[title={Typische Fallen}]
\begin{enumerate}[nosep]
\item Zählung ab \textbf{0}: \texttt{[1]} ist das \emph{zweite} Element.
\item \texttt{firstChild} kann ein unsichtbarer \textbf{Whitespace-Textknoten} sein.
\item Attribute sind \textbf{keine} Kindknoten.
\item \texttt{nodeName} von Elementen ist in \textbf{GROSSBUCHSTABEN}; \texttt{nodeValue} von Elementen ist \texttt{null}.
\item Attribut (wie im HTML) $\neq$ Property (vom Browser bearbeitet), z.\,B.\ \texttt{href} mit Schrägstrich.
\item \texttt{prompt} und \texttt{input.value} liefern \textbf{Text}; \lstinline|"4" + 1| ist \lstinline|"41"|.
\item \lstinline|Number("")| und \lstinline|Number(null)| sind \lstinline|0|, nicht NaN.
\item \lstinline|NaN === NaN| ist \lstinline|false| -- \lstinline|isNaN| verwenden.
\item Listener-Funktion ohne \texttt{()} übergeben.
\item Klicks \textbf{blubbern nach oben} -- ein Klick auf ein inneres Element löst auch Listener außen aus.
\item \lstinline|-1 % 3| ist in JS \lstinline|-1| -- vorher die Länge addieren.
\item Skript im \texttt{<head>} ohne \texttt{defer} findet die Elemente im Body noch nicht.
\end{enumerate}
\end{merkbox}

% =====================================================================
\section{Aufgabe 1: Moodle-Quiz "`JavaScript Concepts"' (0,1 P)}\label{sec:a1}
% =====================================================================

\textbf{In einfachen Worten:} Du machst in Moodle ein kurzes Online-Quiz über JavaScript-Grundbegriffe und gibst
einen Screenshot deines Punktestands ab.

\textbf{Vorgehen:}
\begin{enumerate}[nosep]
\item Quiz in Moodle öffnen (Link steht in \texttt{aufgabe1/LIESMICH.txt}).
\item Vorher die Grundlagen aus Abschnitt~\ref{sec:grundlagen} durchlesen -- vor allem Variablen, Datentypen, Funktionen.
\item Quiz machen, Screenshot der Punkte speichern (z.\,B.\ \texttt{quiz-ergebnis.png}) und mit abgeben.
\end{enumerate}

\begin{tutor}
\textbf{Wie viele Punkte hattest du, und welche Frage war schwierig?} -- Hier ehrlich antworten; das Quiz ist zum
Selbsttesten da.
\end{tutor}

% =====================================================================
\section{Aufgabe 2: Zahlenratespiel (0,1 P)}\label{sec:a2}
% =====================================================================

\subsection{Was bedeutet die Aufgabe in einfachen Worten?}

Der Computer denkt sich eine Zahl von 1 bis 10 aus. Du hast \textbf{3 Versuche}, sie zu erraten. Nach jedem Tipp
sagt er dir "`zu hoch"' oder "`zu niedrig"' und zeigt deine bisherigen Tipps. Errätst du sie, gratuliert er dir;
schaffst du es nicht, verrät er die Zahl. Danach fragt er, ob du nochmal spielen willst. Alles läuft über die
kleinen Browser-Fenster \texttt{prompt}, \texttt{alert} und \texttt{confirm}; die HTML-Seite selbst ist leer.

\subsection{Welche Methode?}

Die Angabe gibt den Aufbau schon vor: \textbf{drei Teile}, die ineinandergreifen.
\begin{enumerate}[nosep]
\item \textbf{Spielschleife} (oben im Skript): \texttt{while}-Schleife "`solange der Spieler will"'. Pro Runde:
      Zufallszahl ziehen, \texttt{playNumberGame} aufrufen, mit \texttt{confirm} fragen.
\item \textbf{\texttt{playNumberGame(secretNumber)}}: eine Runde mit bis zu 3 Versuchen (\texttt{for}-Schleife).
\item \textbf{\texttt{getUserInput()}}: fragt eine Zahl ab und prüft sie mit \texttt{try}/\texttt{catch}.
\end{enumerate}

\begin{analogie}
Das ist wie ein \textbf{Spieleabend}: Die Spielleiterin (Spielschleife) fragt nach jeder Runde "`Noch eine?"'.
Eine Runde (\texttt{playNumberGame}) hat 3 Züge. Bei jedem Zug nimmt ein Schiedsrichter (\texttt{getUserInput})
deinen Tipp entgegen und prüft, ob er überhaupt erlaubt ist.
\end{analogie}

\subsection{Code Schritt für Schritt}

\textbf{Die HTML-Datei} \texttt{aufgabe2/index.html}:
\begin{htmlcode}
<!DOCTYPE html>
<html lang="en">
<head>
    <meta charset="UTF-8">
    <title>Number-Guessing Game</title>
    <script src="game.js" defer></script>
</head>
<body></body>
</html>
\end{htmlcode}
Der Body ist leer, so wie es die Angabe verlangt -- das ganze Spiel passiert in den Dialogfenstern.
Zeile~6 bindet \texttt{game.js} ein.

\subsubsection*{Teil 1: Konstanten und Spielschleife (Zeilen 1--10)}
\begin{code}
const MAX_ATTEMPTS = 3;
// changes after every round, therefore let instead of const
let keepPlaying = true;

while (keepPlaying) {
    let randomNumber = Math.floor(Math.random() * 10) + 1;
    playNumberGame(randomNumber);
    keepPlaying = confirm("Do you want to play again?");
}
alert("Thank you for playing!");
\end{code}

\begin{itemize}
\item \textbf{Zeile 1:} Die erlaubte Versuchsanzahl 3 ändert sich nie, deshalb \texttt{const}. Großbuchstaben sind
      eine übliche Schreibweise für solche festen Werte. \emph{Warum überhaupt eine Konstante statt überall "`3"'?}
      Will man später 5 Versuche, ändert man nur diese eine Zeile.
\item \textbf{Zeile 3:} Die Angabe verlangt eine "`Konstante"' für die Entscheidung "`weiterspielen?"'. Das
      widerspricht sich aber: Diese Entscheidung \emph{ändert} sich nach jeder Runde (Zeile 8), und eine
      \texttt{const} kann man nicht neu zuweisen (das gäbe einen TypeError). Deshalb \texttt{let}. Der Startwert
      \lstinline|true| sorgt dafür, dass mindestens eine Runde gespielt wird.
\item \textbf{Zeile 5:} Die Schleife läuft, solange \texttt{keepPlaying} wahr ist. Die Bedingung wird \emph{vor}
      jeder Runde geprüft.
\item \textbf{Zeile 6:} Zufallszahl 1--10, genau wie in Abschnitt~\ref{sec:zufall} vorgerechnet. \texttt{let} im
      Block: Jede Runde bekommt ihre eigene, frische Variable.
\item \textbf{Zeile 7:} Eine Runde spielen. Die Funktion steht erst weiter unten -- das geht wegen \emph{Hoisting}
      (Funktionsdeklarationen sind von Anfang an bekannt).
\item \textbf{Zeile 8:} \texttt{confirm} liefert \lstinline|true| (OK) oder \lstinline|false| (Abbrechen). Das Ergebnis
      wird direkt in \texttt{keepPlaying} gespeichert. Bei \lstinline|false| ist die Schleifenbedingung beim nächsten
      Prüfen falsch, und die Schleife endet.
\item \textbf{Zeile 10:} Erst nach der Schleife, also wenn der Spieler aufhört: Dankesnachricht.
\end{itemize}

\subsubsection*{Teil 2: Eine Runde -- \texttt{playNumberGame} (Zeilen 12--37)}
\begin{code}[firstnumber=12]
function playNumberGame(secretNumber) {
    let guesses = [];

    for (let attempt = 1; attempt <= MAX_ATTEMPTS; attempt++) {
        let guess = getUserInput();
        guesses.push(guess);

        if (typeof guess !== "number") {
            // guess contains the error message
            alert(guess);
        } else if (guess === secretNumber) {
            alert("Correct! You needed " + attempt + " attempt(s).");
            return true;
        } else {
            if (guess > secretNumber) {
                alert("Too high! Attempts made: " + attempt);
            } else {
                alert("Too low! Attempts made: " + attempt);
            }
            alert("Your previous guesses: " + guesses);
        }
    }

    alert("No attempts left. The correct number was " + secretNumber + ".");
    return false;
}
\end{code}

\begin{itemize}
\item \textbf{Zeile 13:} Leeres Array für alle Tipps dieser Runde. Es steht \emph{in} der Funktion, damit jede neue
      Runde mit einer leeren Liste beginnt.
\item \textbf{Zeile 15:} Schleife über die erlaubten Versuche: \texttt{attempt} läuft 1, 2, 3. Ab 1 statt 0 zu zählen
      hat einen Grund: So kann man \texttt{attempt} direkt als "`Versuch Nummer \dots"' ausgeben.
\item \textbf{Zeilen 16--17:} Tipp holen und in die Liste schreiben (\texttt{push}), wie verlangt.
\item \textbf{Zeile 19:} \texttt{getUserInput} liefert bei ungültiger Eingabe den Fehler\emph{text} (Typ string)
      statt einer Zahl. \lstinline|typeof guess !== "number"| erkennt genau diesen Fall. Dann zeigt Zeile 21 den
      Fehlertext an. Der Versuch ist trotzdem verbraucht, denn die Schleife zählt einfach weiter.
\item \textbf{Zeilen 22--24:} Richtig geraten: Glückwunsch inklusive Anzahl der Versuche, dann \lstinline|return true|.
      \textbf{Warum \texttt{return}?} Es beendet die Funktion -- und damit auch die \texttt{for}-Schleife -- sofort.
      Ohne \texttt{return} würde man trotz richtiger Antwort weiter gefragt.
\item \textbf{Zeilen 25--32:} Falsch geraten: "`zu hoch"' oder "`zu niedrig"' mit Versuchszahl, danach die bisherigen
      Tipps. \lstinline|"..." + guesses| wandelt das Array automatisch in Text mit Kommas um (z.\,B.\ \texttt{3,5}).
\item \textbf{Zeilen 35--36:} Wird nur erreicht, wenn alle 3 Versuche verbraucht sind (sonst hätte
      \lstinline|return true| die Funktion vorher beendet). Dann die richtige Zahl verraten und \lstinline|return false|.
\end{itemize}

\subsubsection*{Teil 3: Eingabe prüfen -- \texttt{getUserInput} (Zeilen 39--49)}
\begin{code}[firstnumber=39]
function getUserInput() {
    try {
        let guess = Number(prompt("Enter a number between 1 and 10:"));
        if (isNaN(guess) || guess < 1 || guess > 10) {
            throw new Error("Invalid input. Please enter a number between 1 and 10.");
        }
        return guess;
    } catch (error) {
        return error.message;
    }
}
\end{code}

\begin{itemize}
\item \textbf{Zeile 41:} \texttt{prompt} liefert Text (oder \texttt{null} bei "`Abbrechen"'). \texttt{Number(...)}
      macht eine Zahl daraus, damit man vergleichen kann.
\item \textbf{Zeile 42:} Drei Gründe für "`ungültig"', verbunden mit "`oder"': keine Zahl (\texttt{isNaN}),
      kleiner als 1, größer als 10. Eines reicht.
\item \textbf{Zeile 43:} Bei ungültiger Eingabe einen Fehler werfen -- \texttt{return guess} in Zeile 45 wird
      dann übersprungen, und es geht direkt in den \texttt{catch}-Block.
\item \textbf{Zeile 47:} Der \texttt{catch}-Block gibt den Fehlertext zurück. So bekommt \texttt{playNumberGame}
      in jedem Fall etwas zurück: entweder eine gültige Zahl oder den Fehlertext.
\end{itemize}

\textbf{Welche Eingaben sind gültig?}
\begin{center}
\begin{tabular}{@{}llll@{}}
\toprule
Eingabe & \texttt{prompt} liefert & \texttt{Number(...)} & Ergebnis\\
\midrule
\texttt{5}          & \lstinline|"5"|   & \lstinline|5|   & gültig, Zahl 5\\
\texttt{abc}        & \lstinline|"abc"| & \lstinline|NaN| & ungültig (\texttt{isNaN})\\
(leer, OK)          & \lstinline|""|    & \lstinline|0|   & ungültig ($0 < 1$)\\
"`Abbrechen"'       & \lstinline|null|  & \lstinline|0|   & ungültig ($0 < 1$)\\
\texttt{11}         & \lstinline|"11"|  & \lstinline|11|  & ungültig ($11 > 10$)\\
\bottomrule
\end{tabular}
\end{center}

\subsection{Ergebnis: zwei getestete Spielverläufe}

\begin{ergebnis}
\textbf{Lauf 1} -- Geheimzahl 5, Eingaben \texttt{3}, dann \texttt{5}, danach "`nicht nochmal"':
\begin{enumerate}[nosep]
\item Versuch 1: \texttt{3} ist gültig, $3 < 5$ $\to$ \texttt{Too low! Attempts made: 1}
\item \texttt{Your previous guesses: 3}
\item Versuch 2: \texttt{5} $=$ Geheimzahl $\to$ \texttt{Correct! You needed 2 attempt(s).} (\texttt{return true})
\item \texttt{confirm} $\to$ Abbrechen $\to$ \texttt{Thank you for playing!}
\end{enumerate}

\textbf{Lauf 2} -- Geheimzahl 5, Eingaben \texttt{9}, \texttt{abc}, Abbrechen:
\begin{enumerate}[nosep]
\item Versuch 1: $9 > 5$ $\to$ \texttt{Too high! Attempts made: 1}, dann \texttt{Your previous guesses: 9}
\item Versuch 2: \texttt{abc} $\to$ NaN $\to$ Fehler $\to$ \texttt{Invalid input. Please enter a number between 1 and 10.}
\item Versuch 3: Abbrechen $\to$ \texttt{null} $\to$ 0 $\to$ derselbe Fehlertext
\item Keine Versuche mehr $\to$ \texttt{No attempts left. The correct number was 5.} (\texttt{return false})
\item \texttt{confirm} $\to$ Abbrechen $\to$ \texttt{Thank you for playing!}
\end{enumerate}
\end{ergebnis}

\begin{tutor}
\begin{enumerate}[leftmargin=*]
\item \textbf{Warum ist \texttt{keepPlaying} \texttt{let} und nicht \texttt{const}, obwohl die Angabe "`constant"' sagt?}\\
  Weil sich der Wert nach jeder Runde durch \texttt{confirm} ändert. Eine \texttt{const} darf man nicht neu
  zuweisen, das gäbe einen TypeError. \texttt{MAX\_ATTEMPTS} dagegen ändert sich nie und ist \texttt{const}.
\item \textbf{Wie entsteht genau eine Zahl von 1 bis 10?}\\
  \texttt{Math.random()} liefert $[0,1)$, mal 10 ergibt $[0,10)$, \texttt{Math.floor} rundet ab auf $0\dots9$,
  $+1$ verschiebt auf $1\dots10$. Die 10 entsteht z.\,B.\ aus $0{,}95$: $9{,}5 \to 9 \to 10$.
\item \textbf{Was passiert, wenn man "`Abbrechen"' drückt?}\\
  \texttt{prompt} liefert \texttt{null}, \texttt{Number(null)} ist 0, und $0 < 1$ ist ungültig. Also wird der
  Fehler geworfen, der Fehlertext angezeigt, und der Versuch zählt als verbraucht.
\item \textbf{Wozu \texttt{try}/\texttt{catch}, wenn man auch einfach \texttt{if} nehmen könnte?}\\
  Die Angabe verlangt es. Der Vorteil: Der Fehlerfall wird an \emph{einer} Stelle (\texttt{catch}) behandelt, und
  \texttt{throw} springt sofort dorthin, ohne dass man ein \texttt{return guess} versehentlich erreicht.
\item \textbf{Warum \texttt{typeof guess !== \textquotedbl{}number\textquotedbl{}} und nicht \texttt{isNaN(guess)}?}\\
  \texttt{getUserInput} gibt im Fehlerfall einen \emph{Text} zurück, nie NaN. Ein Text hat den Typ
  \texttt{\textquotedbl{}string\textquotedbl{}}, eine gültige Zahl den Typ \texttt{\textquotedbl{}number\textquotedbl{}}. So unterscheidet man die beiden Fälle eindeutig.
  (Nebenbei: Weil die Angabe jeden Rückgabewert ins Array schreiben lässt, landen auch Fehlertexte in
  \texttt{guesses} und würden bei einem späteren "`previous guesses"' mit angezeigt.)
\end{enumerate}
\end{tutor}

% =====================================================================
\section{Aufgabe 3: DOM-Operationen (0,2 P)}\label{sec:a3}
% =====================================================================

\subsection{Was bedeutet die Aufgabe in einfachen Worten?}

Es gibt eine kleine Seite \texttt{ex3.html} mit drei Links, die je ein Flaggenbild enthalten. Du sollst
\begin{enumerate}[nosep]
\item[(1)] den \textbf{DOM-Baum} dieser Seite zeichnen (wie auf Folie 3-29),
\item[(2)] ein gegebenes Skript einbinden, die Seite öffnen und \textbf{jede der 16 Konsolenausgaben} genau erklären,
\item[(3)] erklären, was zwei weitere gegebene Zeilen bewirken.
\end{enumerate}
Hier wird nichts programmiert -- es geht ums \emph{Verstehen}, wie der Browser die Seite intern speichert.

\subsection{Die Seite und das Skript}

\texttt{aufgabe3/ex3.html} (Bild-Adressen gekürzt; Zeile 6 wurde zum Einbinden ergänzt):
\begin{htmlcode}
<!DOCTYPE html>
<html lang="en">
<head>
    <meta charset="UTF-8">
    <title>DOM Operations</title>
    <script src="dom.js" defer></script>
</head>
<body>
<a href=https://en.wikipedia.org target="_new"><img
        src="https://.../uk_..._british_flag-512.png" alt="British Flag"></a>
BeforeParent
<a href="https://de.wikipedia.org" target="_new"><p>Text</p>
    Before <img src="https://.../austria-national-world-512.png" alt="Austrian Flag">
    After
</a>
AfterParent
<a href="https://ja.wikipedia.org" target="_new">
    <img src="https://.../Flags_Flag_Country_National_Symbol_Japan-512.png" alt="Japanese Flag">
</a>
</body>
</html>
\end{htmlcode}

\textbf{Warum \texttt{defer}?} Ohne \texttt{defer} liefe \texttt{dom.js} schon im \texttt{<head>}, bevor die Links
überhaupt gelesen sind. \lstinline|getElementsByTagName("a")[1]| wäre dann \lstinline|undefined|, und die erste
\texttt{console.log}-Zeile würde mit einem Fehler (TypeError) abbrechen.

\subsection{Teil 3.1: Der DOM-Baum}

\textbf{Vorgehen beim Zeichnen} (Strategie für jede HTML-Datei):
\begin{merkbox}[title={Strategie: DOM-Baum zeichnen}]
\begin{enumerate}[nosep]
\item Ganz oben \texttt{\#document}, darunter \texttt{<!DOCTYPE>} und \texttt{html}.
\item Jedes Tag wird ein Elementknoten; alles, was \emph{zwischen} seinem Start- und End-Tag steht, sind seine Kinder.
\item Jeder Text zwischen zwei Tags wird ein Textknoten -- \textbf{auch reiner Whitespace} (Zeilenumbruch, Einrückung).
\item Attribute schreibt man \emph{neben} das Element, nicht als Kind.
\item Kinder in der Reihenfolge von links nach rechts bzw.\ oben nach unten, wie sie im Quelltext stehen.
\end{enumerate}
\end{merkbox}

Der folgende Baum wurde im Browser für die \emph{originale} Datei (ohne die eingefügte \texttt{script}-Zeile)
nachgemessen. Blau = Elementknoten, orange = Textknoten, grün in eckigen Klammern = Attribute (keine Kinder!).
\nl\ = Zeilenumbruch, \ws\ = Leerzeichen.

\begin{center}
\begin{tikzpicture}[
  x=0.62cm, y=0.5cm,
  every node/.style={anchor=west, inner sep=2pt, font=\ttfamily\footnotesize},
  el/.style={draw=blue!60, fill=blue!8, rounded corners=1pt},
  tx/.style={draw=orange!70, fill=orange!8, rounded corners=1pt},
  dt/.style={draw=gray!60, fill=gray!8, rounded corners=1pt},
  attr/.style={text=green!45!black, font=\ttfamily\scriptsize},
  hl/.style={draw=red!70, very thick, fill=blue!8, rounded corners=1pt},
  ann/.style={text=red!70!black, font=\sffamily\scriptsize}]
\newcommand{\edge}[2]{\draw[gray] ([xshift=5pt]#1.south west) |- (#2.west);}
% rows
\node[dt] (doc)  at (0,0)   {\#document};
\node[dt] (dtp)  at (1,-1)  {<!DOCTYPE html>};
\node[el] (html) at (1,-2)  {html};  \node[attr] at (html.east) {[lang=\qq{en}]};
\node[el] (head) at (2,-3)  {head};
\node[tx] (h1)   at (3,-4)  {\#text \qq{\nl\ws\ws\ws\ws}};
\node[el] (meta) at (3,-5)  {meta}; \node[attr] at (meta.east) {[charset=\qq{UTF-8}]};
\node[tx] (h2)   at (3,-6)  {\#text \qq{\nl\ws\ws\ws\ws}};
\node[el] (tit)  at (3,-7)  {title};
\node[tx] (tt)   at (4,-8)  {\#text \qq{DOM Operations}};
\node[tx] (hb)   at (2,-9)  {\#text \qq{\nl}};
\node[el] (body) at (2,-10) {body};
\node[tx] (b1)   at (3,-11) {\#text \qq{\nl}};
\node[el] (a1)   at (3,-12) {a}; \node[attr] at (a1.east) {[href=\qq{https://en.wikipedia.org}, target=\qq{\_new}]};
\node[el] (i1)   at (4,-13) {img}; \node[attr] at (i1.east) {[src=\qq{...}, alt=\qq{British Flag}]};
\node[tx] (b2)   at (3,-14) {\#text \qq{\nl BeforeParent\nl}};
\node[hl] (a2)   at (3,-15) {a}; \node[attr] (a2at) at (a2.east) {[href=\qq{https://de.wikipedia.org}, target=\qq{\_new}]};
\node[ann] at (a2at.east) {$\leftarrow$ link (Index 1)};
\node[el] (p)    at (4,-16) {p};
\node[tx] (pt)   at (5,-17) {\#text \qq{Text}};
\node[tx] (c2)   at (4,-18) {\#text \qq{\nl\ws\ws\ws\ws Before\ws}};
\node[el] (i2)   at (4,-19) {img}; \node[attr] at (i2.east) {[src=\qq{...}, alt=\qq{Austrian Flag}]};
\node[tx] (c3)   at (4,-20) {\#text \qq{\nl\ws\ws\ws\ws After\nl}};
\node[tx] (b3)   at (3,-21) {\#text \qq{\nl AfterParent\nl}};
\node[el] (a3)   at (3,-22) {a}; \node[attr] (a3at) at (a3.east) {[href=\qq{https://ja.wikipedia.org}, target=\qq{\_new}]};
\node[ann] at (a3at.east) {$\leftarrow$ item(2)};
\node[tx] (d1)   at (4,-23) {\#text \qq{\nl\ws\ws\ws\ws}};
\node[el] (i3)   at (4,-24) {img}; \node[attr] at (i3.east) {[src=\qq{...}, alt=\qq{Japanese Flag}]};
\node[tx] (d2)   at (4,-25) {\#text \qq{\nl}};
\node[tx] (b4)   at (3,-26) {\#text \qq{\nl\nl}};
% edges
\edge{doc}{dtp} \edge{doc}{html}
\edge{html}{head} \edge{html}{hb} \edge{html}{body}
\edge{head}{h1} \edge{head}{meta} \edge{head}{h2} \edge{head}{tit}
\edge{tit}{tt}
\edge{body}{b1} \edge{body}{a1} \edge{body}{b2} \edge{body}{a2} \edge{body}{b3} \edge{body}{a3} \edge{body}{b4}
\edge{a1}{i1}
\edge{a2}{p} \edge{a2}{c2} \edge{a2}{i2} \edge{a2}{c3}
\edge{p}{pt}
\edge{a3}{d1} \edge{a3}{i3} \edge{a3}{d2}
\end{tikzpicture}
\end{center}

\textbf{Was man am Baum erklären können sollte:}
\begin{itemize}
\item \textbf{Der zweite Link hat vier Kinder}: \texttt{<p>}, den Textknoten "`\nl\ws\ws\ws\ws Before\ws"', das
      \texttt{<img>} und den Textknoten "`\nl\ws\ws\ws\ws After\nl"'. Vor dem \texttt{<p>} gibt es \emph{keinen}
      Whitespace-Knoten, weil \texttt{<p>} im Quelltext direkt nach \texttt{<a ...>} steht.
\item \textbf{Der erste Link hat nur ein Kind} (\texttt{img}), weil der Zeilenumbruch im Quelltext \emph{innerhalb}
      des \texttt{<img ...>}-Tags liegt (zwischen \texttt{<img} und \texttt{src}). Dort ist er nur Trennzeichen
      zwischen Attributen und wird kein Textknoten.
\item \textbf{Der dritte Link} hat dagegen vor und nach dem Bild Whitespace-Textknoten, weil \texttt{<img>} dort
      in einer eigenen, eingerückten Zeile steht.
\item \textbf{Text zwischen den Links} ("`BeforeParent"', "`AfterParent"') sind Textknoten direkt im
      \texttt{body}, samt Zeilenumbrüchen davor und danach.
\item \textbf{Besonderheiten des Parsers:} Der Zeilenumbruch zwischen \texttt{</head>} und \texttt{<body>} landet als
      Textknoten im \texttt{html}-Element. Die Zeilenumbrüche nach \texttt{</body>} landen laut HTML-Regeln
      trotzdem im \texttt{body} (letzter Knoten "`\nl\nl"').
\item \textbf{Attribute} (\texttt{href}, \texttt{target}, \texttt{src}, \texttt{alt}, \texttt{lang}, \texttt{charset})
      hängen am Element, sind aber \emph{keine} Kindknoten.
\end{itemize}

\subsection{Teil 3.2: Die 16 Konsolenausgaben}

Das Skript \texttt{aufgabe3/dom.js} (Teil 3.2):
\begin{code}
// Exercise 3.2
var link = document.getElementsByTagName("a")[1];
console.log(link.attributes["href"].nodeValue);
console.log(link.getAttributeNode("href").nodeValue);
console.log(link.href);
console.log(link.firstChild);
console.log(link.attributes[0].nodeType);
console.log(link.firstChild.nodeType);
console.log(link.parentNode.nodeName);
console.log(link.lastChild.parentNode.nodeName);
console.log(link.attributes.length);
console.log(link.innerHTML);
console.log(link.outerHTML);
console.log(link.innerText);
console.log(link.outerText);
console.log(link.textContent);
console.log(link);
console.log(link.nodeValue);
\end{code}

\textbf{Zeile 2 zuerst:} \lstinline|getElementsByTagName("a")| liefert alle drei Links. \lstinline|[1]| ist wegen der
Zählung ab 0 der \textbf{zweite} Link, also der deutsche (\texttt{de.wikipedia.org}) mit \texttt{<p>Text</p>}
darin. Alle folgenden Zeilen beziehen sich auf genau diesen Link (im Baum rot umrandet).

In den Ausgaben unten steht \texttt{\textbackslash n} für einen Zeilenumbruch. Die Konsole zeigt an dieser Stelle
eine echte neue Zeile an; zur besseren Lesbarkeit sind die Texte hier in Anführungszeichen und mit
\texttt{\textbackslash n} geschrieben.

\subsubsection*{Ausgaben 1--3: Attribut gegen Property}

\textbf{1.} \lstinline|link.attributes["href"].nodeValue|
\begin{ausgabe}
https://de.wikipedia.org
\end{ausgabe}
\lstinline|attributes| ist die Liste der Attributknoten; \lstinline|["href"]| holt den Knoten mit dem Namen
\texttt{href}. Der \texttt{nodeValue} eines Attributknotens ist sein Wert -- \textbf{genau wie im HTML geschrieben},
also ohne Schrägstrich am Ende.

\textbf{2.} \lstinline|link.getAttributeNode("href").nodeValue|
\begin{ausgabe}
https://de.wikipedia.org
\end{ausgabe}
Derselbe Attributknoten, nur mit einer anderen Methode geholt. Daher dasselbe Ergebnis.

\textbf{3.} \lstinline|link.href|
\begin{ausgabe}
https://de.wikipedia.org/
\end{ausgabe}
Das ist die \textbf{DOM-Property}, nicht das Attribut. Der Browser liefert die aufgelöste, vereinheitlichte (normalisierte)
absolute Adresse; dabei ergänzt er den Pfad "`/"'. Deshalb steht am Ende ein Schrägstrich, der im HTML nicht steht.

\subsubsection*{Ausgaben 4--9: Navigation im Baum und Knotentypen}

\textbf{4.} \lstinline|link.firstChild|
\begin{ausgabe}
<p>Text</p>
\end{ausgabe}
Das erste Kind des Links ist das \texttt{<p>}-Element. Die Konsole zeigt es als aufklappbares Element an (kein Text,
sondern ein Objekt). Kein Whitespace-Knoten davor, weil \texttt{<p>} im Quelltext direkt auf \texttt{<a ...>} folgt.

\textbf{5.} \lstinline|link.attributes[0].nodeType|
\begin{ausgabe}
2
\end{ausgabe}
\lstinline|attributes[0]| ist das erste Attribut in der Reihenfolge des HTML, also \texttt{href}. Attributknoten haben
den Typ \textbf{2} (\texttt{ATTRIBUTE\_NODE}).

\textbf{6.} \lstinline|link.firstChild.nodeType|
\begin{ausgabe}
1
\end{ausgabe}
Das erste Kind ist das \texttt{<p>} (siehe Ausgabe 4), ein Element. Elemente haben den Typ \textbf{1}
(\texttt{ELEMENT\_NODE}).

\textbf{7.} \lstinline|link.parentNode.nodeName|
\begin{ausgabe}
BODY
\end{ausgabe}
Der Elternknoten des Links ist \texttt{<body>} (siehe Baum). \texttt{nodeName} von HTML-Elementen ist immer in
\textbf{Großbuchstaben}, obwohl im Quelltext \texttt{body} klein geschrieben ist.

\textbf{8.} \lstinline|link.lastChild.parentNode.nodeName|
\begin{ausgabe}
A
\end{ausgabe}
Schritt für Schritt: \lstinline|link.lastChild| ist das letzte Kind des Links, der Textknoten
"`\nl\ws\ws\ws\ws After\nl"'. Dessen \lstinline|parentNode| ist wieder der Link selbst. Also landet man beim Link,
und sein \texttt{nodeName} ist \texttt{\textquotedbl{}A\textquotedbl{}}. ("`Vom Vater zum letzten Kind und zurück zum Vater."')

\textbf{9.} \lstinline|link.attributes.length|
\begin{ausgabe}
2
\end{ausgabe}
Der Link hat genau zwei Attribute: \texttt{href} und \texttt{target}.

\subsubsection*{Ausgaben 10--14: Inhalte als Text}

\textbf{10.} \lstinline|link.innerHTML| (Bild-Adresse gekürzt)
\begin{ausgabe}
"<p>Text</p>\n    Before <img src="https://cdn3.iconfinder.com/.../austria-national-world-512.png" alt="Austrian Flag">\n    After\n"
\end{ausgabe}
Der \textbf{HTML-Code innerhalb} des \texttt{<a>} als Text, ohne das \texttt{<a>}-Tag selbst. Der Whitespace aus dem
Quelltext (Zeilenumbrüche, 4 Leerzeichen Einrückung) bleibt erhalten.

\textbf{11.} \lstinline|link.outerHTML| (Bild-Adresse gekürzt)
\begin{ausgabe}
"<a href="https://de.wikipedia.org" target="_new"><p>Text</p>\n    Before <img src="https://cdn3.iconfinder.com/.../austria-national-world-512.png" alt="Austrian Flag">\n    After\n</a>"
\end{ausgabe}
Wie \texttt{innerHTML}, aber \textbf{inklusive} Start- und End-Tag des Links selbst.

\textbf{12.} \lstinline|link.innerText|
\begin{ausgabe}
"Text\n\nBefore  After "
\end{ausgabe}
Das ist der Text, \textbf{wie man ihn auf dem Bildschirm sieht}. Schritt für Schritt:
\begin{enumerate}[nosep]
\item \texttt{<p>Text</p>} liefert "`Text"'. Weil \texttt{<p>} ein \textbf{Blockelement} (Absatz) ist, setzt
      \texttt{innerText} danach \textbf{zwei} Zeilenumbrüche -- so wie zwischen zwei Absätzen eine Leerzeile ist.
\item Der Textknoten "`\nl\ws\ws\ws\ws Before\ws"': Beim Anzeigen werden Zeilenumbruch und Einrückung am Zeilenanfang
      verschluckt (Whitespace wird zusammengefasst). Übrig bleibt "`Before\ws"' mit \emph{einem} Leerzeichen am Ende.
\item Das \texttt{<img>} hat keinen Text, es trägt nichts bei.
\item Der Textknoten "`\nl\ws\ws\ws\ws After\nl"' wird zu "`\ws After\ws"': Die Folge aus Zeilenumbruch und vier
      Leerzeichen schrumpft zu einem Leerzeichen, der Zeilenumbruch am Ende ebenfalls.
\item Ergebnis: Zwischen "`Before"' und "`After"' stehen \textbf{zwei} Leerzeichen (je eines links und rechts vom
      Bild). Am Ende bleibt ein Leerzeichen, weil der Link ein Inline-Element ist und in derselben Zeile direkt
      "`AfterParent"' folgt -- der Abstand wird also tatsächlich angezeigt.
\end{enumerate}

\textbf{13.} \lstinline|link.outerText|
\begin{ausgabe}
"Text\n\nBefore  After "
\end{ausgabe}
Beim \textbf{Lesen} identisch mit \texttt{innerText}. Der Unterschied gibt es nur beim \emph{Setzen}:
\lstinline|link.outerText = "x"| würde den ganzen Link durch den Text ersetzen, \lstinline|link.innerText = "x"| nur seinen Inhalt.

\textbf{14.} \lstinline|link.textContent|
\begin{ausgabe}
"Text\n    Before \n    After\n"
\end{ausgabe}
Alle Textknoten \textbf{roh} hintereinander, ohne auf das Aussehen zu achten: "`Text"' (aus dem \texttt{<p>}) $+$
"`\nl\ws\ws\ws\ws Before\ws"' $+$ (Bild: nichts) $+$ "`\nl\ws\ws\ws\ws After\nl"'. Kein doppelter Zeilenumbruch
nach "`Text"', weil \texttt{textContent} nicht weiß, dass \texttt{<p>} ein Absatz ist. Der \texttt{alt}-Text
"`Austrian Flag"' fehlt, weil er ein Attribut ist und kein Textknoten.

\subsubsection*{Ausgaben 15--16: das Element selbst}

\textbf{15.} \lstinline|link|
\begin{ausgabe}
<a href="https://de.wikipedia.org" target="_new">...</a>
\end{ausgabe}
Das Element-Objekt selbst. Die Konsole zeigt es aufklappbar an; klappt man es auf, sieht man die Kinder.

\textbf{16.} \lstinline|link.nodeValue|
\begin{ausgabe}
null
\end{ausgabe}
\textbf{Elementknoten haben keinen \texttt{nodeValue}} -- der ist nur bei Text-, Attribut- und Kommentarknoten belegt.
Den Text eines Elements bekommt man über \texttt{textContent} oder \texttt{innerText}.

\begin{ergebnis}[title={Ergebnis 3.2 -- alle 16 Ausgaben auf einen Blick}]
\small
\begin{tabular}{@{}rL{6.2cm}L{7.2cm}@{}}
\toprule
Nr. & Ausdruck & Ausgabe\\
\midrule
1  & \lstinline|link.attributes["href"].nodeValue| & \texttt{https://de.wikipedia.org}\\
2  & \lstinline|link.getAttributeNode("href")...| & \texttt{https://de.wikipedia.org}\\
3  & \lstinline|link.href| & \texttt{https://de.wikipedia.org/} (mit \texttt{/})\\
4  & \lstinline|link.firstChild| & Element \texttt{<p>Text</p>}\\
5  & \lstinline|link.attributes[0].nodeType| & \texttt{2} (Attribut)\\
6  & \lstinline|link.firstChild.nodeType| & \texttt{1} (Element)\\
7  & \lstinline|link.parentNode.nodeName| & \texttt{BODY}\\
8  & \lstinline|link.lastChild.parentNode.nodeName| & \texttt{A}\\
9  & \lstinline|link.attributes.length| & \texttt{2}\\
10 & \lstinline|link.innerHTML| & HTML im Link (\texttt{<p>}\dots\texttt{After\textbackslash n})\\
11 & \lstinline|link.outerHTML| & HTML inkl.\ \texttt{<a ...>}\dots\texttt{</a>}\\
12 & \lstinline|link.innerText| & \texttt{\textquotedbl{}Text\textbackslash n\textbackslash nBefore\ws\ws After\ws\textquotedbl{}}\\
13 & \lstinline|link.outerText| & wie 12\\
14 & \lstinline|link.textContent| & \texttt{\textquotedbl{}Text\textbackslash n\ws\ws\ws\ws Before\ws\textbackslash n\ws\ws\ws\ws After\textbackslash n\textquotedbl{}}\\
15 & \lstinline|link| & das \texttt{<a>}-Element (aufklappbar)\\
16 & \lstinline|link.nodeValue| & \texttt{null}\\
\bottomrule
\end{tabular}
\end{ergebnis}

\subsection{Teil 3.3: \texttt{removeAttribute(\textquotedbl{}target\textquotedbl{})}}

\begin{code}[firstnumber=20]
// Exercise 3.3
var link = document.getElementsByTagName("a").item(2);
link.removeAttribute("target");
\end{code}

\textbf{Schritt für Schritt:}
\begin{enumerate}
\item \lstinline|.item(2)| ist dasselbe wie \lstinline|[2]|: wegen Zählung ab 0 der \textbf{dritte} Link, also der
      japanische (\texttt{ja.wikipedia.org}).
\item \lstinline|removeAttribute("target")| löscht bei \emph{diesem} Link das Attribut \texttt{target=\textquotedbl{}\_new\textquotedbl{}}.
      Adresse (\texttt{href}) und Bild bleiben unverändert, ebenso die beiden anderen Links (die behalten
      \texttt{target=\textquotedbl{}\_new\textquotedbl{}}). Im Browser nachgeprüft: Die \texttt{target}-Werte der drei Links sind danach
      \texttt{[\textquotedbl{}\_new\textquotedbl{}, \textquotedbl{}\_new\textquotedbl{}, null]}.
\item \textbf{Was ändert sich für den Benutzer?} Ohne \texttt{target} öffnet ein Link im \textbf{aktuellen Tab}
      (Standardverhalten; es gibt auch kein \texttt{<base target>} auf der Seite, das etwas anderes festlegen würde).
      Klickt man also auf die japanische Flagge, wird die eigene Seite durch Wikipedia ersetzt.
\end{enumerate}

\textbf{Was bedeutet \texttt{target=\textquotedbl{}\_new\textquotedbl{}} überhaupt -- und warum ist das nicht dasselbe wie \texttt{\_blank}?}
\begin{itemize}
\item \texttt{target=\textquotedbl{}\_blank\textquotedbl{}} ist ein \textbf{reserviertes Spezialwort}: Jeder Klick öffnet einen \emph{neuen},
      namenlosen Tab.
\item \texttt{\_new} ist \textbf{kein} Spezialwort, sondern einfach ein \textbf{Name} für ein Fenster/einen Tab.
      Beim ersten Klick gibt es noch keinen Tab mit diesem Namen, also öffnet der Browser einen neuen und nennt ihn
      "`\_new"'. Jeder weitere Klick auf einen Link mit \texttt{target=\textquotedbl{}\_new\textquotedbl{}} lädt dann in \textbf{genau diesen}
      Tab, statt einen weiteren zu öffnen.
\end{itemize}

\begin{analogie}
\texttt{\_blank} ist wie "`Leg jedes Mal ein \textbf{neues} Blatt Papier hin."' \texttt{\_new} ist wie "`Schreib
auf das Blatt mit der Aufschrift \emph{\_new}; gibt es noch keines, leg eins hin und beschrifte es."'
Nach \texttt{removeAttribute} schreibt der japanische Link einfach auf das Blatt, das man gerade vor sich hat.
\end{analogie}

\textbf{Warum darf \texttt{var link} zweimal stehen?} Bei \texttt{var} ist es erlaubt, dieselbe Variable nochmals
zu deklarieren -- sie wird einfach überschrieben. Mit \texttt{let} wäre das ein SyntaxError, und das ganze Skript
würde gar nicht laufen.

\begin{tutor}
\begin{enumerate}[leftmargin=*]
\item \textbf{Warum ist \texttt{link.firstChild} hier das \texttt{<p>} und kein Textknoten?}\\
  Weil im Quelltext \texttt{<p>} direkt nach \texttt{<a ...>} steht, ohne Leerzeichen oder Zeilenumbruch
  dazwischen. Es entsteht kein Whitespace-Knoten.
\item \textbf{Warum unterscheiden sich Ausgabe 1 und 3?}\\
  1 liest das \emph{Attribut} (Text wie im HTML), 3 die \emph{Property} (vom Browser normalisierte absolute
  URL mit Pfad "`/"').
\item \textbf{Was ist der Unterschied zwischen \texttt{innerText} und \texttt{textContent}?}\\
  \texttt{textContent} liefert alle Textknoten roh mit dem originalen Whitespace; \texttt{innerText} liefert den
  sichtbaren Text, Whitespace wie beim Anzeigen zusammengefasst, mit Zeilenumbrüchen nach Blockelementen.
\item \textbf{Warum ist \texttt{link.nodeValue} null?}\\
  Elemente haben keinen eigenen Wert, nur Kinder. \texttt{nodeValue} gibt es bei Text-, Attribut- und
  Kommentarknoten.
\item \textbf{Was ändert sich durch 3.3?}\\
  Nur beim dritten Link verschwindet \texttt{target}; er öffnet nun im aktuellen Tab statt im Tab namens
  "`\_new"'. \texttt{\_new} ist kein Spezialwort wie \texttt{\_blank}, sondern ein Fenstername.
\end{enumerate}
\end{tutor}

% =====================================================================
\section{Aufgabe 4: Glaspreis (0,2 P)}\label{sec:a4}
% =====================================================================

\subsection{Was bedeutet die Aufgabe in einfachen Worten?}

Es gibt ein Bestellformular für Glasscheiben (\texttt{glassOrder.html}). Man wählt Farbe, Dicke, Breite und Höhe.
Die Seite soll \textbf{sofort den Preis anzeigen}, sobald man etwas ändert. Die Bedingung: Die HTML-Datei darf
\textbf{nicht verändert} werden -- alles muss das JavaScript in \texttt{calculate.js} erledigen:
\begin{enumerate}[nosep]
\item den Absatz löschen, der die Überschrift wiederholt (\texttt{<p>Your Glass Order</p>}),
\item ein schreibgeschütztes Preisfeld mit Beschriftung "`price"' ins Formular einbauen,
\item den Preis berechnen: $\text{Preis} = \dfrac{\text{Breite} \cdot \text{Höhe} \cdot \text{Dicke}}{32}$ bei farbigem
      Glas, $\dfrac{\dots}{50}$ bei einem Grauton; alles in cm,
\item nur einen Preis anzeigen, wenn alle nötigen Eingaben da und gültig sind, sonst einen Bindestrich "`-"',
\item das bei \emph{jeder} Änderung automatisch neu machen.
\end{enumerate}

\begin{analogie}
Das ist wie die \textbf{Waage beim Obststand}, die sofort den Preis anzeigt: Sobald du etwas auf die Waage legst oder
wegnimmst (Eingabe ändern), rechnet sie neu (\texttt{update}). Liegt nichts Sinnvolles drauf, zeigt sie "`-"'.
\end{analogie}

\subsection{Welche Methode?}
\begin{itemize}[nosep]
\item Löschen und Einfügen: DOM-Methoden \texttt{querySelector}, \texttt{remove}, \texttt{createElement},
      \texttt{append} (Abschnitt~\ref{sec:finden}).
\item Gültigkeit: \texttt{checkValidity()} (Abschnitt~\ref{sec:formulare}).
\item Automatisch neu rechnen: Event Listener für \texttt{input} und \texttt{change} (Abschnitt~\ref{sec:events}).
\item Grauton erkennen: Hex-Farbe in drei Teile zerlegen und vergleichen.
\end{itemize}

\subsection{Vorab: eine Beispielrechnung von Hand}

Bevor wir den Code ansehen, rechnen wir ein Beispiel mit \emph{anderen} Zahlen als in den Tests komplett durch:
Breite 50\,cm, Höhe 20\,cm, Dicke 8\,mm, Farbe \texttt{\#00ff00} (Grün).

\textbf{Schritt 1 -- Grauton oder bunt?} \texttt{\#00ff00} zerlegt: Rot \texttt{00}, Grün \texttt{ff}, Blau \texttt{00}.
\texttt{00} $\neq$ \texttt{ff}, also nicht alle gleich $\Rightarrow$ \textbf{bunt}, Teiler 32.

\textbf{Schritt 2 -- Einheiten:} Breite und Höhe sind laut Formular schon in cm. Die Dicke ist in mm
(die Radio-Buttons heißen 4\,mm, 6\,mm, 8\,mm, 12\,mm). Weil $1\,\text{cm} = 10\,\text{mm}$, gilt
$8\,\text{mm} = 8 / 10 = 0{,}8\,\text{cm}$.

\textbf{Schritt 3 -- rechnen:}
\begin{align*}
\text{Preis} &= \frac{b \cdot h \cdot d}{32} && \text{bunt, also durch 32}\\
&= \frac{50 \cdot 20 \cdot 0{,}8}{32} && \text{Werte einsetzen, alles in cm}\\
&= \frac{800}{32} && 50 \cdot 20 = 1000,\quad 1000 \cdot 0{,}8 = 800\\
&= 25 && 32 \cdot 25 = 800\\
&\to \texttt{25.00} && \text{Anzeige mit genau 2 Nachkommastellen}
\end{align*}

\textbf{Grauton-Beispiel:} \texttt{\#808080} $\to$ Rot \texttt{80}, Grün \texttt{80}, Blau \texttt{80}. Alle drei gleich
$\Rightarrow$ Grauton, Teiler 50. Genauso \texttt{\#000000} (Schwarz) und \texttt{\#ffffff} (Weiß): Auch das sind
Graustufen, weil $R = G = B$.

\subsection{Code Schritt für Schritt}

\subsubsection*{Block 1: Elemente merken (Zeilen 1--4)}
\begin{code}
const widthInput = document.getElementById("width");
const heightInput = document.getElementById("height");
const colorInput = document.getElementById("color");
const thicknessInputs = document.querySelectorAll('input[name="thickness"]');
\end{code}
Die Eingabefelder werden \emph{einmal} gesucht und in Variablen gespeichert, weil man sie später oft braucht.
Breite, Höhe und Farbe haben eine \texttt{id}, also \texttt{getElementById}. Die vier Radio-Buttons haben
keine gemeinsame \texttt{id}, aber denselben \texttt{name=\textquotedbl{}thickness\textquotedbl{}} -- \texttt{querySelectorAll} mit
Attribut-Selektor holt alle vier als Liste.

\subsubsection*{Block 2: doppelten Absatz löschen (Zeile 7)}
\begin{code}[firstnumber=6]
// delete the paragraph that repeats the heading
document.querySelector("h1 + p").remove();
\end{code}
Im HTML steht \lstinline[language=HTML]|<h1>Your Glass Order</h1>| und direkt darunter
\lstinline[language=HTML]|<p>Your Glass Order</p>|. Der Absatz hat keine \texttt{id}. Der Selektor \texttt{h1 + p}
bedeutet "`das \texttt{<p>}, das \emph{direkt nach} einem \texttt{<h1>} kommt"' -- genau dieser Absatz.
\texttt{remove()} entfernt ihn aus dem DOM, und er verschwindet von der Seite. Die HTML-\emph{Datei} bleibt dabei
unverändert, nur der Baum im Browser ändert sich.

\subsubsection*{Block 3: Preisfeld einbauen (Zeilen 9--20)}
\begin{code}[firstnumber=9]
// add read-only price field with label at the end of the glass specification
const priceInput = document.createElement("input");
priceInput.type = "text";
priceInput.id = "price";
priceInput.name = "price";
priceInput.readOnly = true;

const priceLabel = document.createElement("label");
priceLabel.htmlFor = "price";
priceLabel.textContent = "price";

widthInput.parentElement.append(priceInput, priceLabel);
\end{code}
\begin{itemize}
\item \textbf{Zeilen 10--14:} Ein neues \texttt{<input>} bauen (noch unsichtbar). Typ Text, \texttt{id} und
      \texttt{name} "`price"'. \lstinline|readOnly = true| macht es \textbf{schreibgeschützt}: Der Benutzer kann den
      Preis nicht selbst eintippen, nur das Skript kann ihn setzen -- sonst könnte man sich einen beliebigen
      Preis hineinschreiben.
\item \textbf{Zeilen 16--18:} Ein \texttt{<label>} mit Text "`price"'. \lstinline|htmlFor = "price"| verbindet es mit
      dem Feld (klickt man auf die Beschriftung, springt der Cursor ins Feld). Die Property heißt \texttt{htmlFor},
      weil \texttt{for} in JS ein reserviertes Wort ist.
\item \textbf{Zeile 20:} \lstinline|widthInput.parentElement| ist das \texttt{<fieldset>} "`Glass Specification"', in dem
      das Breitenfeld steckt. \texttt{append} hängt Feld und Label als letzte Kinder an. Das ist genau die Stelle
      nach der Legende "`Price"' und dem Kommentar \texttt{<!--Input the new element here via JavaScript-->}, denn
      beide sind die letzten Knoten im Fieldset. (Auch ein Kommentar ist ein Knoten im DOM, Typ 8 -- \texttt{append}
      fügt danach ein.) Warum über \texttt{widthInput}? Das Fieldset hat keine \texttt{id}; über das bekannte
      Breitenfeld findet man es eindeutig.
\end{itemize}

\subsubsection*{Block 4: Preis berechnen -- \texttt{calculate} (Zeilen 22--34)}
\begin{code}[firstnumber=22]
// width, height and thickness in cm
function calculate(width, height, thickness) {
    // color value has the form #rrggbb
    const color = colorInput.value.toLowerCase();
    const r = color.substring(1, 3);
    const g = color.substring(3, 5);
    const b = color.substring(5, 7);

    if (r === g && g === b) {
        return width * height * thickness / 50; // shade of grey
    }
    return width * height * thickness / 32; // colored
}
\end{code}
\begin{itemize}
\item \textbf{Zeile 25:} Farbwert holen, z.\,B.\ \texttt{\textquotedbl{}\#808080\textquotedbl{}}. \texttt{toLowerCase()} wandelt in Kleinbuchstaben um. Das ist eine Absicherung, damit
      \texttt{\textquotedbl{}FF\textquotedbl{}} und \texttt{\textquotedbl{}ff\textquotedbl{}} sicher als gleich gelten (das Farbfeld liefert ohnehin Kleinbuchstaben).
\item \textbf{Zeilen 26--28:} \lstinline|substring(start, ende)| schneidet ein Stück Text heraus, \emph{von} Position
      \texttt{start} \emph{bis vor} Position \texttt{ende}. Positionen zählen ab 0:
      \begin{center}
      \begin{tabular}{@{}lccccccc@{}}
      \toprule
      Position & 0 & 1 & 2 & 3 & 4 & 5 & 6\\
      \midrule
      Zeichen  & \texttt{\#} & \texttt{8} & \texttt{0} & \texttt{8} & \texttt{0} & \texttt{8} & \texttt{0}\\
      Teil     & -- & \multicolumn{2}{c}{\texttt{r} = (1,3)} & \multicolumn{2}{c}{\texttt{g} = (3,5)} & \multicolumn{2}{c}{\texttt{b} = (5,7)}\\
      \bottomrule
      \end{tabular}
      \end{center}
      Das \texttt{\#} an Position 0 wird übersprungen; jede Farbe besteht aus genau zwei Zeichen.
\item \textbf{Zeile 30:} Sind alle drei Teile gleich ($R = G = B$), ist es ein Grauton. Man vergleicht die Texte
      direkt -- eine Umrechnung in Dezimalzahlen ist unnötig, weil gleiche Hex-Texte gleiche Zahlen bedeuten.
      Zwei Vergleiche reichen: Wenn \texttt{r} gleich \texttt{g} ist und \texttt{g} gleich \texttt{b}, ist auch \texttt{r} gleich \texttt{b}.
\item \textbf{Zeilen 31 und 33:} Formel mit Teiler 50 (Grau) bzw.\ 32 (bunt). Das erste \texttt{return} beendet die
      Funktion; deshalb braucht man kein \texttt{else}.
\end{itemize}

\subsubsection*{Block 5: Prüfen und anzeigen -- \texttt{update} (Zeilen 36--50)}
\begin{code}[firstnumber=36]
function update() {
    // a required radio group is only valid if one radio is checked
    const valid = widthInput.checkValidity() &&
        heightInput.checkValidity() &&
        thicknessInputs[0].checkValidity();

    if (!valid) {
        priceInput.value = "-";
        return;
    }

    const thickness = document.querySelector('input[name="thickness"]:checked').value / 10; // mm -> cm
    const price = calculate(Number(widthInput.value), Number(heightInput.value), thickness);
    priceInput.value = price.toFixed(2);
}
\end{code}
\begin{itemize}
\item \textbf{Zeilen 38--40:} Gültig nur, wenn Breite, Höhe \emph{und} Dicke gültig sind. \texttt{checkValidity()}
      prüft die Regeln aus dem HTML: \texttt{required}, \texttt{min=\textquotedbl{}0\textquotedbl{}}, \texttt{max=\textquotedbl{}200\textquotedbl{}}, \texttt{step=\textquotedbl{}1\textquotedbl{}}.
      Leer, 300 oder 2.5 sind also ungültig.
\item \textbf{Warum \texttt{thicknessInputs[0]}?} \texttt{required} steht am ersten Radio (4\,mm) und gilt damit für die
      ganze Gruppe. Ein Radio einer required-Gruppe ist genau dann gültig, wenn \emph{irgendein} Radio der Gruppe
      ausgewählt ist. Es reicht also, eines davon zu fragen; \texttt{[0]} ist einfach das erste.
\item \textbf{Warum nur diese drei Felder?} Für den Preis zählen nur die preisrelevanten Pflichtfelder. Name,
      Adresse usw.\ beeinflussen den Preis nicht. Die Farbe hat immer einen Wert (Standard \texttt{\#000000}),
      sie kann nicht fehlen.
\item \textbf{Zeilen 42--45:} Ungültig $\Rightarrow$ Bindestrich anzeigen und mit \texttt{return} sofort aufhören, damit
      nicht mit fehlenden Werten gerechnet wird.
\item \textbf{Zeile 47:} Den ausgewählten Radio-Button finden (\texttt{:checked}), seinen Wert holen (Text, z.\,B.\
      \texttt{\textquotedbl{}4\textquotedbl{}}) und durch 10 teilen: mm $\to$ cm. Die Division wandelt den Text automatisch in eine Zahl um:
      \lstinline|"4" / 10| ergibt \lstinline|0.4|.
\item \textbf{Zeile 48:} Breite und Höhe sind Texte; \texttt{Number(...)} wandelt sie ausdrücklich um, dann rechnet
      \texttt{calculate}.
\item \textbf{Zeile 49:} \lstinline|toFixed(2)| macht daraus einen Text mit genau zwei Nachkommastellen
      (40 $\to$ \texttt{\textquotedbl{}40.00\textquotedbl{}}, 62.5 $\to$ \texttt{\textquotedbl{}62.50\textquotedbl{}}), wie man es von Preisen kennt.
\end{itemize}

\subsubsection*{Block 6: Event Listener (Zeilen 52--58)}
\begin{code}[firstnumber=52]
widthInput.addEventListener("input", update);
heightInput.addEventListener("input", update);
colorInput.addEventListener("input", update);
for (const radio of thicknessInputs) {
    radio.addEventListener("change", update);
}
update();
\end{code}
\begin{itemize}
\item \textbf{Zeilen 52--54:} \texttt{input} feuert bei jeder Änderung, auch bei jedem einzelnen Tastendruck in
      Breite/Höhe und beim Ziehen im Farbwähler. So ist der Preis immer aktuell.
\item \textbf{Zeilen 55--57:} Jeder der vier Radio-Buttons bekommt einen \texttt{change}-Listener; der feuert, wenn
      dieser Radio ausgewählt wird. Die Schleife spart, viermal dieselbe Zeile zu schreiben.
\item \textbf{Zeile 58:} \texttt{update()} einmal sofort aufrufen. Am Anfang ist noch nichts eingegeben, also zeigt
      das Preisfeld gleich "`-"' statt leer zu sein.
\end{itemize}

\subsection{Ergebnis: die getesteten Preise}

Alle Tests mit Breite 100\,cm und Höhe 50\,cm.

\textbf{Test 1:} Dicke 4\,mm, Farbe \texttt{\#000000} (Schwarz $\Rightarrow$ $R=G=B=$ \texttt{00} $\Rightarrow$ Grauton)
\begin{align*}
\text{Preis} &= \frac{b \cdot h \cdot d}{50} && \text{Grauton, also durch 50}\\
&= \frac{100 \cdot 50 \cdot 0{,}4}{50} && 4\,\text{mm} = 4/10 = 0{,}4\,\text{cm}\\
&= \frac{2000}{50} && 100\cdot 50 = 5000,\quad 5000 \cdot 0{,}4 = 2000\\
&= 40 && \to \text{Anzeige } \texttt{40.00}
\end{align*}

\textbf{Test 2:} Dicke 4\,mm, Farbe \texttt{\#ff0000} (Rot: \texttt{ff} $\neq$ \texttt{00} $\Rightarrow$ bunt)
\begin{align*}
\text{Preis} &= \frac{100 \cdot 50 \cdot 0{,}4}{32} && \text{bunt, also durch 32}\\
&= \frac{2000}{32} && \text{Zähler wie oben}\\
&= 62{,}5 && 32 \cdot 62{,}5 = 2000\\
&\to \texttt{62.50} && \text{zwei Nachkommastellen}
\end{align*}

\textbf{Test 3:} Dicke 12\,mm, Farbe \texttt{\#808080} (\texttt{80} $=$ \texttt{80} $=$ \texttt{80} $\Rightarrow$ Grauton)
\begin{align*}
\text{Preis} &= \frac{100 \cdot 50 \cdot 1{,}2}{50} && 12\,\text{mm} = 1{,}2\,\text{cm},\ \text{Grauton}\\
&= \frac{6000}{50} && 5000 \cdot 1{,}2 = 6000\\
&= 120 && \to \text{Anzeige } \texttt{120.00}
\end{align*}

\begin{ergebnis}
\begin{tabular}{@{}lll@{}}
\toprule
Eingabe & Preisfeld & Begründung\\
\midrule
$100 \times 50$, 4\,mm, \texttt{\#000000} & \texttt{40.00}  & Grauton, $/50$\\
$100 \times 50$, 4\,mm, \texttt{\#ff0000} & \texttt{62.50}  & bunt, $/32$\\
$100 \times 50$, 12\,mm, \texttt{\#808080} & \texttt{120.00} & Grauton, $/50$\\
Breite 300                                & \texttt{-}      & $300 > \texttt{max}=200$\\
Breite leer                               & \texttt{-}      & \texttt{required} verletzt\\
nur Breite gesetzt                        & \texttt{-}      & Höhe und Dicke fehlen\\
\bottomrule
\end{tabular}

Die HTML-Datei ist unverändert (bytegleich geprüft, also Zeichen für Zeichen identisch mit dem Original).
Achtung: Das Formular hat eine \texttt{mailto}-Aktion, d.\,h.\ beim Absenden würde der Browser ein E-Mail-Programm
öffnen -- beim Testen deshalb nicht auf "`Absenden"' klicken.
\end{ergebnis}

\begin{tutor}
\begin{enumerate}[leftmargin=*]
\item \textbf{Wie erkennst du einen Grauton?}\\
  Das Farbfeld liefert \texttt{\#rrggbb}. Mit \texttt{substring} hole ich die drei Zweiergruppen. Sind alle gleich
  ($R=G=B$), ist es Grau, z.\,B.\ \texttt{\#808080}. Schwarz und Weiß zählen auch dazu.
\item \textbf{Warum teilst du die Dicke durch 10?}\\
  Die Formel braucht alles in cm. Breite und Höhe sind schon in cm, die Dicke ist in mm; $1\,\text{cm} = 10\,\text{mm}$.
\item \textbf{Warum prüfst du nur \texttt{thicknessInputs[0]}?}\\
  \texttt{required} gilt für die ganze Radio-Gruppe. Jedes Radio der Gruppe ist gültig, sobald irgendeines
  ausgewählt ist, also reicht eines.
\item \textbf{Wie hast du das Preisfeld eingefügt, ohne das HTML zu ändern?}\\
  Mit \texttt{createElement} gebaut, \texttt{readOnly} gesetzt und mit \texttt{append} ans Ende des Fieldsets
  "`Glass Specification"' gehängt (gefunden über \texttt{widthInput.parentElement}).
\item \textbf{Warum \texttt{input} und nicht \texttt{change} bei Breite/Höhe?}\\
  \texttt{input} feuert bei jeder Änderung sofort (jeder Tastendruck); \texttt{change} erst, wenn man das Feld
  verlässt. Die Angabe will Neuberechnung "`bei jeder Wertänderung"'.
\end{enumerate}
\end{tutor}

% =====================================================================
\section{Aufgabe 5: Event Listener und Target (0,2 P)}\label{sec:a5}
% =====================================================================

\subsection{Was bedeutet die Aufgabe in einfachen Worten?}

Auf der Seite ist ein blauer Kasten (\texttt{<div id=\textquotedbl{}outerArea\textquotedbl{}>}) und darin ein Button "`Click me!"'. Sowohl am
Button als auch am Kasten hängt ein Klick-Listener, der Infos in die Konsole schreibt. Du sollst
\begin{enumerate}[nosep]
\item vorhersagen, was beim Klick auf den Button ausgegeben wird,
\item es ausprobieren und vergleichen,
\item erklären, was \texttt{target}, \texttt{currentTarget} und \texttt{this} sind,
\item den Code so \emph{ergänzen}, dass beim Button-Klick \textbf{nur} die Infos vom Button erscheinen
      (ohne \texttt{ex5.html} zu ändern und ohne den bestehenden Code zu ändern).
\end{enumerate}

\subsection{Der gegebene Code}

\texttt{ex5.html} (gekürzt auf den Body):
\begin{htmlcode}[numbers=none]
<div id="outerArea">
    <button id="button">Click me!</button>
</div>
\end{htmlcode}

\texttt{ex5.js}, Zeilen 1--13 (Original):
\begin{code}
document.querySelector('#button').addEventListener('click',function (event){
    console.log("inner:");
    console.log("target: "+event.target);
    console.log("currentTarget: "+event.currentTarget);
    console.log("this: "+this);
});

document.querySelector('#outerArea').addEventListener('click',function (event){
        console.log("outer:");
        console.log("target: " + event.target);
        console.log("currentTarget: " + event.currentTarget);
        console.log("this: " + this);
});
\end{code}

\subsection{5.1 Vorhersage -- Schritt für Schritt}

\begin{merkbox}[title={Strategie: Event-Ausgaben vorhersagen}]
\begin{enumerate}[nosep]
\item Welches Element wird \emph{tatsächlich} angeklickt? $\to$ das ist \texttt{target}, für alle Listener gleich.
\item Weg nach oben aufschreiben (Bubbling): Element, Elternteil, \dots, \texttt{document}.
\item Auf diesem Weg die Elemente mit Listener suchen -- in dieser Reihenfolge laufen sie.
\item Bei jedem Listener: \texttt{currentTarget} $=$ \texttt{this} $=$ das Element, an dem er hängt.
\item Achtung, wie ausgegeben wird: Bei \texttt{\textquotedbl{}text\textquotedbl{} + element} erscheint \texttt{[object ...Element]}.
\end{enumerate}
\end{merkbox}

Angewendet:
\begin{enumerate}
\item Angeklickt wird der \texttt{<button>} $\Rightarrow$ \texttt{target} ist in beiden Listenern der Button.
\item Weg nach oben: \texttt{button} $\to$ \texttt{div\#outerArea} $\to$ \texttt{body} $\to$ \texttt{html} $\to$ \texttt{document}.
\item Listener gibt es am Button ("`inner"') und am div ("`outer"'). Also läuft zuerst "`inner"', dann "`outer"'.
\item Im inner-Listener ist \texttt{currentTarget} der Button, im outer-Listener das div.
\item Die Werte werden mit \texttt{+} an einen Text gehängt. Dabei wird das Element in Text umgewandelt, und
      daraus wird nur seine Typbezeichnung: \texttt{[object HTMLButtonElement]} bzw.\ \texttt{[object HTMLDivElement]}.
\end{enumerate}

\begin{center}
\begin{tikzpicture}[
  box/.style={draw, rounded corners, minimum width=3.4cm, minimum height=0.75cm, font=\ttfamily\small},
  lis/.style={font=\sffamily\footnotesize, text=red!70!black, anchor=west},
  node distance=0.45cm]
\node[box, fill=gray!10]  (doc)  {document};
\node[box, fill=gray!10, below=of doc]  (html) {html};
\node[box, fill=gray!10, below=of html] (body) {body};
\node[box, fill=blue!15, below=of body] (div)  {div\#outerArea};
\node[box, fill=orange!25, below=of div] (btn) {button\#button};
\foreach \a/\b in {btn/div, div/body, body/html, html/doc}
  \draw[-{Stealth[length=2.5mm]}, thick, blue!60!black] (\a.north) -- (\b.south);
\node[lis] at ([xshift=0.4cm]btn.east) {1. Listener "`inner"' läuft (+ in 5.4: stopPropagation)};
\node[lis] at ([xshift=0.4cm]div.east) {2. Listener "`outer"' läuft};
\node[lis, text=gray] at ([xshift=0.4cm]body.east) {kein Listener};
\node[font=\sffamily\small, rotate=90, text=blue!60!black] at ([xshift=-0.9cm]div.west) {Bubbling: nach oben};
\node[font=\sffamily\scriptsize, below=0.15cm of btn] {Klick passiert hier $\Rightarrow$ \texttt{target}};
\end{tikzpicture}
\end{center}

\subsection{5.2 Tatsächliche Ausgabe}

Die Vorhersage stimmt genau mit der Ausgabe im Browser überein:
\begin{ausgabe}
inner:
target: [object HTMLButtonElement]
currentTarget: [object HTMLButtonElement]
this: [object HTMLButtonElement]
outer:
target: [object HTMLButtonElement]
currentTarget: [object HTMLDivElement]
this: [object HTMLDivElement]
\end{ausgabe}

\textbf{Gegenprobe:} Klickt man direkt auf die blaue Fläche (nicht auf den Button), läuft nur der outer-Listener,
und \texttt{target}, \texttt{currentTarget} und \texttt{this} sind alle drei das div
(\texttt{[object HTMLDivElement]}), denn das div selbst wurde angeklickt.

\subsection{5.3 \texttt{target}, \texttt{currentTarget} und \texttt{this}}

\begin{center}
\begin{tabular}{@{}lL{7.2cm}L{3.2cm}@{}}
\toprule
 & Bedeutung & bei Button-Klick im outer-Listener\\
\midrule
\texttt{event.target} & Element, auf dem das Ereignis \textbf{ausgelöst} wurde (das tatsächlich angeklickte,
                        innerste Element). Bleibt während des ganzen Bubblings gleich. & Button\\
\texttt{event.currentTarget} & Element, an dem der \textbf{gerade laufende Listener hängt}. Ändert sich beim
                        Hochblubbern. & div\\
\texttt{this} & In einer normalen \texttt{function} als Listener dasselbe wie \texttt{currentTarget}. (Bei einer
                Pfeilfunktion wäre es \emph{nicht} das Element.) & div\\
\bottomrule
\end{tabular}
\end{center}

\subsection{5.4 Nur die Button-Infos ausgeben}

\textbf{Idee:} Das Problem ist das Bubbling: Der Klick steigt vom Button zum div auf und löst dort "`outer"' aus.
Wenn man das Aufsteigen \emph{am Button} stoppt, kommt der Klick nie beim div an.

Ergänzung in \texttt{ex5.js} ab Zeile 14 (der Originalcode bleibt unverändert):
\begin{code}[firstnumber=15]
// Exercise 5.4: stop the click on the button from bubbling up to outerArea
document.querySelector('#button').addEventListener('click', function (event) {
    event.stopPropagation();
});
\end{code}

\begin{itemize}
\item Ein \emph{zusätzlicher} Listener am Button. Man darf an einem Element beliebig viele Listener anhängen;
      sie laufen in der Reihenfolge, in der sie angehängt wurden.
\item \lstinline|event.stopPropagation()| verhindert, dass das Event zum div (und weiter nach oben) gelangt.
\item \textbf{Warum läuft "`inner"' trotzdem?} \texttt{stopPropagation} stoppt nur den Weg zu \emph{anderen}
      Elementen, nicht weitere Listener am selben Element. Beide Listener am Button laufen also noch.
\item \textbf{Warum ein neuer Listener statt Änderung des alten?} Die Angabe sagt "`nur Code hinzufügen"'. Der neue
      Listener erfüllt das, ohne eine einzige alte Zeile anzufassen.
\end{itemize}

\begin{ergebnis}
Button-Klick $\Rightarrow$ nur noch die vier Zeilen \texttt{inner:}, \texttt{target: [object HTMLButtonElement]},
\texttt{currentTarget: [object HTMLButtonElement]}, \texttt{this: [object HTMLButtonElement]}.\\
Klick auf die blaue Fläche $\Rightarrow$ weiterhin \texttt{outer:} (dort wird ja nichts gestoppt).
\end{ergebnis}

\textbf{Alternative} (nicht gewählt): Im outer-Listener am Anfang \lstinline|if (event.target !== event.currentTarget) return;|
-- "`nur reagieren, wenn das div selbst angeklickt wurde"'. Dafür müsste man aber den bestehenden Code ändern,
was die Angabe nicht erlaubt.

\begin{tutor}
\begin{enumerate}[leftmargin=*]
\item \textbf{Warum wird beim Button-Klick auch "`outer"' ausgegeben?}\\
  Wegen Bubbling: Das click-Event steigt vom Button über das div, body, html bis document auf, und jeder Listener
  auf dem Weg läuft -- auch der am div.
\item \textbf{Was ist der Unterschied zwischen \texttt{target} und \texttt{currentTarget}?}\\
  \texttt{target} ist das angeklickte Element (bleibt gleich), \texttt{currentTarget} das Element, an dem der gerade
  laufende Listener hängt (wandert mit).
\item \textbf{Was ist \texttt{this}?}\\
  In einer normalen \texttt{function} als Listener dasselbe wie \texttt{currentTarget}. In einer Pfeilfunktion nicht.
\item \textbf{Warum steht da \texttt{[object HTMLButtonElement]} und nicht der Button selbst?}\\
  Wegen der String-Verkettung mit \texttt{+}: Das Element wird in Text umgewandelt, und das ergibt nur seinen Typnamen.
  Mit \texttt{console.log(event.target)} ohne \texttt{+} sähe man das Element aufklappbar.
\item \textbf{Läuft der inner-Listener nach deiner Änderung noch?}\\
  Ja. \texttt{stopPropagation} verhindert nur, dass das Event zu anderen Elementen weiterwandert; alle Listener am
  Button selbst laufen.
\end{enumerate}
\end{tutor}

% =====================================================================
\section{Aufgabe 6: Bildergalerie (0,2 P)}\label{sec:a6}
% =====================================================================

\subsection{Was bedeutet die Aufgabe in einfachen Worten?}

Es gibt eine Galerie aus der Vorlesung: Kleine Vorschaubilder (\emph{Thumbnails}); klickt man eines an, erscheint
es groß über die ganze Seite (Vollbildmodus); ein Klick irgendwo schließt das große Bild wieder. Du sollst
\begin{enumerate}[nosep]
\item verstehen, wie HTML und \texttt{gallery.js} aufgebaut sind,
\item im Vollbildmodus zwei Buttons "`previous"' und "`next"' ergänzen: "`next"' zeigt das nächste Bild, nach dem
      letzten wieder das erste; "`previous"' das vorige, vor dem ersten das letzte,
\item das so machen, dass es mit \emph{beliebig vielen} Bildern im gleichen Format funktioniert,
\item dabei nur Code \emph{hinzufügen}, nichts am bestehenden Code ändern.
\end{enumerate}

\begin{analogie}
Das ist wie ein \textbf{Karussell} oder ein Diaprojektor mit rundem Magazin: Nach dem letzten Dia kommt wieder das
erste, und wenn man beim ersten zurückdreht, landet man beim letzten. Rechnerisch ist das die Modulo-Uhr aus
Abschnitt~\ref{sec:modulo}.
\end{analogie}

\subsection{Verstehen: HTML-Struktur}
\begin{htmlcode}[firstnumber=9]
    <div class="gallery">
       <img src="hk_small.JPG" data-large="hk.JPG" alt="View of Hong Kong from the peak">
     <img src="street_small.JPG" data-large="street.JPG" alt="Hong Kong central district life">
     <img src="taichung_small.JPG" data-large="taichung.JPG" alt="A teahouse in Taiwan">
  </div>
     <div id="displayDIV">
         <img class="bigIMG" src="" alt="">
         <p class="desc">My Caption</p>
         <!-- Add the buttons dynamically with js here -->
     </div>
\end{htmlcode}
\begin{itemize}
\item \texttt{div.gallery}: die Thumbnails. Jedes \texttt{<img>} hat drei wichtige Attribute: \texttt{src} (kleine
      Datei, die man sieht), \texttt{data-large} (Dateiname des großen Bilds) und \texttt{alt} (Beschreibung).
      Attribute, die mit \texttt{data-} beginnen, sind \textbf{eigene Datenattribute}: Man darf sie frei erfinden,
      um Zusatzinfos ans Element zu hängen. In JS liest man sie über \lstinline|element.dataset|:
      \texttt{data-large} wird zu \lstinline|element.dataset.large|.
\item \texttt{div\#displayDIV}: die Vollbild-Ebene mit einem großen Bild (\texttt{.bigIMG}) und einer
      Bildunterschrift (\texttt{.desc}). Laut CSS ist sie zuerst \textbf{unsichtbar} (\texttt{display:none}), deckt
      die ganze Seite ab (\texttt{100vw} $\times$ \texttt{100vh} = 100\,\% der Fensterbreite mal 100\,\% der Fensterhöhe;
      \texttt{position:absolute} mit \texttt{top:0; left:0} legt sie oben links über alles andere) und ist silbern.
\end{itemize}

\subsection{Verstehen: der Originalcode (Zeilen 1--16)}
\begin{code}
function showImg(e) {
   document.querySelector("#displayDIV .bigIMG").src=e.target.dataset.large;
   document.querySelector("#displayDIV .desc").innerHTML=e.target.alt;
   document.getElementById("displayDIV").style.display="block";
}

function hideImg(e) {
   e.currentTarget.style.display="none";
}

let images = document.querySelectorAll('.gallery img');
for (let i=0; i< images.length; i++) {
   images[i].addEventListener("click", showImg);
}
document.getElementById("displayDIV").addEventListener("click", hideImg);
\end{code}
\begin{itemize}
\item \textbf{\texttt{showImg} (Zeilen 1--5):} \texttt{e.target} ist das angeklickte Thumbnail. Das große Bild bekommt
      dessen \texttt{data-large} als Quelle, die Unterschrift dessen \texttt{alt}-Text, und die Vollbild-Ebene wird mit
      \texttt{display=\textquotedbl{}block\textquotedbl{}} sichtbar gemacht.
\item \textbf{\texttt{hideImg} (Zeilen 7--9):} \texttt{e.currentTarget} ist das Element, an dem der Listener hängt, also
      \texttt{displayDIV}. Es wird wieder versteckt. Weil Klicks nach oben blubbern, schließt \emph{jeder} Klick
      irgendwo auf der Ebene (auch auf das große Bild) das Vollbild.
\item \textbf{Zeilen 11--15:} Alle Thumbnails holen und jedem \texttt{showImg} als Klick-Listener geben; der
      Vollbild-Ebene \texttt{hideImg}.
\end{itemize}

\subsection{Lösungsidee}

Was fehlt, ist eine Antwort auf die Frage "`\textbf{Welches} Bild wird gerade gezeigt?"'. Das Original weiß das
nicht -- es liest die Daten nur einmal vom angeklickten Thumbnail. Also:
\begin{enumerate}[nosep]
\item eine Variable \texttt{currentIndex} = Nummer des aktuellen Bilds in der Liste \texttt{images},
\item beim Klick auf ein Thumbnail diese Nummer merken,
\item eine Funktion \texttt{showImgAt(index)}, die das Bild mit dieser Nummer anzeigt -- mit Wrap-around,
\item zwei Buttons, die \texttt{showImgAt(currentIndex - 1)} bzw.\ \texttt{showImgAt(currentIndex + 1)} aufrufen.
\end{enumerate}

\subsection{Code Schritt für Schritt (Ergänzung ab Zeile 17)}

\subsubsection*{Index merken (Zeilen 18--25)}
\begin{code}[firstnumber=17]
// Extension: previous / next buttons
let currentIndex = 0;

// remember which thumbnail was opened
for (let i = 0; i < images.length; i++) {
   images[i].addEventListener("click", function () {
      currentIndex = i;
   });
}
\end{code}
\begin{itemize}
\item \textbf{Zeile 18:} Startwert 0; er wird beim ersten Thumbnail-Klick ohnehin überschrieben.
\item \textbf{Zeilen 21--25:} Jedes Thumbnail bekommt einen \emph{zusätzlichen} Klick-Listener (das Original wird
      nicht verändert, \texttt{showImg} läuft weiterhin zuerst). Dieser merkt sich die Nummer \texttt{i} des Bilds.
\item \textbf{Warum funktioniert \texttt{i} in der Funktion?} Die Funktion läuft erst viel später, beim Klick. Weil
      \texttt{i} mit \texttt{let} angelegt ist, bekommt \textbf{jeder Schleifendurchlauf sein eigenes \texttt{i}}, und
      die Funktion "`erinnert sich"' an genau dieses (das nennt man \emph{Closure}). Mit \texttt{var} gäbe es nur
      ein einziges \texttt{i}, das nach der Schleife 3 ist -- jeder Klick würde dann \texttt{currentIndex = 3} setzen.
\end{itemize}

\begin{analogie}
Closure mit \texttt{let}: Jedes Thumbnail bekommt bei der Schleife einen \textbf{eigenen Zettel} mit seiner Nummer
in die Tasche gesteckt. Mit \texttt{var} gäbe es nur \emph{eine} Tafel an der Wand, auf der am Ende "`3"' steht,
und alle würden beim Klick auf diese Tafel schauen.
\end{analogie}

\subsubsection*{Bild Nummer \texttt{index} anzeigen (Zeilen 27--31)}
\begin{code}[firstnumber=27]
function showImgAt(index) {
   currentIndex = (index + images.length) % images.length; // wraps around at both ends
   document.querySelector("#displayDIV .bigIMG").src = images[currentIndex].dataset.large;
   document.querySelector("#displayDIV .desc").innerHTML = images[currentIndex].alt;
}
\end{code}
\begin{itemize}
\item \textbf{Zeile 28:} Wrap-around (Rechnung siehe unten). Danach liegt \texttt{currentIndex} sicher zwischen 0 und
      \texttt{images.length - 1}.
\item \textbf{Zeilen 29--30:} Genau wie \texttt{showImg} im Original -- nur nimmt man die Daten nicht von
      \texttt{e.target}, sondern vom Thumbnail \texttt{images[currentIndex]}. Das Vollbild ist schon offen, daher
      muss \texttt{display} nicht nochmals gesetzt werden.
\end{itemize}

\textbf{Wrap-around-Rechnung Schritt für Schritt} mit 3 Bildern (Index 0 = \texttt{hk}, 1 = \texttt{street},
2 = \texttt{taichung}), also \texttt{images.length} $= 3$:

\emph{Fall 1: Beim ersten Bild (Index 0) auf "`previous"'} $\Rightarrow$ \texttt{showImgAt(0 - 1)}:
\begin{align*}
\text{index} &= 0 + (-1) = -1 && \text{previous: Schritt } -1\\
(\text{index} + 3) \,\%\, 3 &= (-1 + 3) \,\%\, 3 && \text{Länge addieren, damit nichts Negatives entsteht}\\
&= 2 \,\%\, 3 && 3 \text{ passt 0-mal in } 2\\
&= 2 && \text{Rest } 2 \Rightarrow \text{letztes Bild (taichung)}
\end{align*}

\emph{Fall 2: Beim letzten Bild (Index 2) auf "`next"'} $\Rightarrow$ \texttt{showImgAt(2 + 1)}:
\begin{align*}
\text{index} &= 2 + 1 = 3 && \text{next: Schritt } +1\text{; Index 3 gibt es nicht}\\
(\text{index} + 3) \,\%\, 3 &= (3 + 3) \,\%\, 3 && \text{Länge addieren (schadet hier nicht)}\\
&= 6 \,\%\, 3 && 6 = 2 \cdot 3 + 0\\
&= 0 && \text{Rest } 0 \Rightarrow \text{erstes Bild (hk)}
\end{align*}

\emph{Fall 3: Normalfall, Index 0 auf "`next"'}:
\begin{align*}
(0 + 1 + 3) \,\%\, 3 &= 4 \,\%\, 3 && \text{Länge addieren}\\
&= 1 && 4 = 1 \cdot 3 + 1 \Rightarrow \text{zweites Bild (street)}
\end{align*}

\textbf{Warum das "`$+\,\texttt{images.length}$"'?} Ohne diesen Summanden wäre im Fall 1
\lstinline|-1 % 3| in JavaScript gleich $-1$, und \lstinline|images[-1]| existiert nicht. Das Addieren der Länge
verschiebt um eine volle "`Runde"', was das Ergebnis im Kreis nicht verändert, macht die Zahl aber nicht-negativ.

\subsubsection*{Die Buttons (Zeilen 33--47)}
\begin{code}[firstnumber=33]
function addButton(text, step, side) {
   let button = document.createElement("button");
   button.textContent = text;
   button.style.position = "absolute";
   button.style.top = "50%";
   button.style[side] = "20px";
   button.addEventListener("click", function (e) {
      e.stopPropagation(); // otherwise hideImg on displayDIV closes the picture
      showImgAt(currentIndex + step);
   });
   document.getElementById("displayDIV").appendChild(button);
}

addButton("previous", -1, "left");
addButton("next", 1, "right");
\end{code}
\begin{itemize}
\item \textbf{Warum eine Funktion für beide Buttons?} Beide sind fast gleich; sie unterscheiden sich nur in Text,
      Schrittweite ($-1$ / $+1$) und Seite (links / rechts). Eine Funktion mit drei Parametern vermeidet doppelten Code.
\item \textbf{Zeilen 34--35:} Button bauen und beschriften.
\item \textbf{Zeilen 36--38:} \texttt{position:absolute} positioniert ihn relativ zum nächsten positionierten Vorfahren
      -- das ist \texttt{displayDIV}, das selbst \texttt{position:absolute} hat. \texttt{top: 50\%} setzt ihn auf halbe
      Höhe. \lstinline|button.style[side]| ist die \textbf{Klammer-Schreibweise}: Mit \texttt{side = \textquotedbl{}left\textquotedbl{}} bedeutet
      das dasselbe wie \lstinline|button.style.left|; so kann ein Parameter bestimmen, \emph{welche} Eigenschaft
      gesetzt wird.
\item \textbf{Zeile 40 -- warum \texttt{stopPropagation} nötig ist:} Der Button liegt \emph{innerhalb} von
      \texttt{displayDIV}. Ohne Stopp würde der Klick vom Button zum \texttt{displayDIV} hochblubbern, dort liefe
      \texttt{hideImg}, und das Vollbild würde sofort geschlossen -- man sähe das neue Bild nie. Mit
      \texttt{stopPropagation} kommt der Klick bei \texttt{displayDIV} nicht an. Klicks außerhalb der Buttons
      schließen weiterhin wie im Original.
\item \textbf{Zeile 41:} Neues Bild anzeigen: aktueller Index plus Schritt; den Wrap-around erledigt \texttt{showImgAt}.
\item \textbf{Zeile 43:} Button als letztes Kind in \texttt{displayDIV} einhängen (dort, wo im HTML der Kommentar steht).
      Weil er \emph{in} der Vollbild-Ebene steckt, ist er automatisch nur sichtbar, wenn diese sichtbar ist.
\item \textbf{Zeilen 46--47:} Beide Buttons einmal beim Laden erzeugen.
\end{itemize}

\textbf{Warum funktioniert das mit beliebigem HTML im gleichen Format?} Nirgends steht eine feste Zahl wie "`3"'
oder ein fester Dateiname. Die Liste \texttt{images} wird aus dem Original übernommen
(\lstinline|querySelectorAll('.gallery img')|), und Index-Rechnung sowie Bilddaten hängen nur von
\texttt{images.length}, \texttt{data-large} und \texttt{alt} ab. Mit 5 oder 50 Bildern läuft es genauso.

\begin{ergebnis}
Getestet im Browser:
\begin{itemize}[nosep]
\item erstes Bild öffnen $\to$ next $\to$ \texttt{street.JPG} mit "`Hong Kong central district life"',
\item 2$\times$ previous $\to$ \texttt{taichung.JPG} (Wrap zum letzten),
\item next $\to$ \texttt{hk.JPG} (Wrap zum ersten),
\item Klick außerhalb der Buttons schließt das Vollbild,
\item drittes Bild öffnen $+$ next $\to$ \texttt{hk.JPG}; keine JavaScript-Fehler.
\end{itemize}
\end{ergebnis}

\begin{tutor}
\begin{enumerate}[leftmargin=*]
\item \textbf{Wie funktioniert der Sprung vom letzten zum ersten Bild?}\\
  \lstinline|(index + images.length) % images.length|. Bei 3 Bildern: next von 2 ergibt $(3+3)\,\%\,3 = 0$,
  previous von 0 ergibt $(-1+3)\,\%\,3 = 2$.
\item \textbf{Warum brauchst du \texttt{stopPropagation} beim Button?}\\
  Der Button liegt in \texttt{displayDIV}, an dem \texttt{hideImg} hängt. Ohne Stopp würde der Klick hochblubbern
  und das Vollbild sofort schließen.
\item \textbf{Was machen \texttt{data-large} und \texttt{dataset}?}\\
  \texttt{data-large} ist ein eigenes Datenattribut mit dem Dateinamen des großen Bilds; in JS liest man es mit
  \lstinline|img.dataset.large|.
\item \textbf{Woher weiß dein Code, welches Bild gerade offen ist?}\\
  Ein zusätzlicher Klick-Listener je Thumbnail speichert dessen Index in \texttt{currentIndex}; dank \texttt{let i}
  hat jeder Listener seine eigene Nummer.
\item \textbf{Warum funktioniert es auch mit mehr Bildern?}\\
  Es gibt keine festen Zahlen; alles hängt von \texttt{images.length} und den Attributen der Bilder ab.
\end{enumerate}
\end{tutor}

% =====================================================================
\section{Zusammenfassung}\label{sec:zusammenfassung}
% =====================================================================

\begin{center}
\small
\begin{tabular}{@{}lL{6.0cm}L{7.4cm}@{}}
\toprule
Aufgabe & Kern-Idee & Ergebnis\\
\midrule
1 Quiz & Moodle-Quiz "`JavaScript Concepts"' & selbst machen, Screenshot der Punkte abgeben\\
\addlinespace
2 Ratespiel & \texttt{while}-Spielschleife, \texttt{for} über 3 Versuche, \texttt{try}/\texttt{catch} bei der Eingabe,
  \texttt{Math.floor(Math.random()*10)+1} & z.\,B.\ Geheimzahl 5, Tipps 3, 5 $\to$ "`Too low"', "`Correct! You needed 2 attempt(s)."'\\
\addlinespace
3 DOM & Baum mit Whitespace-Textknoten; Attribut vs.\ Property; 5 Inhalts-Properties & 16 Ausgaben erklärt
  (u.\,a.\ \texttt{href} mit \texttt{/}, \texttt{BODY}, \texttt{2}/\texttt{1}, \texttt{null}); 3.3: dritter Link
  öffnet im aktuellen Tab\\
\addlinespace
4 Glaspreis & \texttt{remove}, \texttt{createElement}/\texttt{append}, \texttt{checkValidity}, Grau bei $R=G=B$,
  Dicke mm$\to$cm & $100\times50$: 4\,mm schwarz \texttt{40.00}, 4\,mm rot \texttt{62.50}, 12\,mm grau
  \texttt{120.00}; ungültig \texttt{-}\\
\addlinespace
5 Events & Bubbling; \texttt{target} bleibt, \texttt{currentTarget}/\texttt{this} wandert & ohne Ergänzung inner $+$
  outer; mit zusätzlichem \texttt{stopPropagation}-Listener nur inner\\
\addlinespace
6 Galerie & \texttt{currentIndex}, \texttt{(i+n)\%n}, Buttons per JS mit \texttt{stopPropagation} &
  previous/next mit Wrap-around, für beliebig viele Bilder\\
\bottomrule
\end{tabular}
\end{center}

\begin{merkbox}[title={Goldene Regeln}]
\begin{enumerate}[nosep]
\item \textbf{Skripte mit \texttt{defer}} einbinden, damit das DOM schon fertig ist.
\item \textbf{\texttt{const}} für Werte, die sich nie ändern; \textbf{\texttt{let}} für alles, was sich ändert;
      \texttt{var} nur kennen, nicht benutzen.
\item \textbf{Eingaben sind Text} -- mit \texttt{Number()} umwandeln und auf \texttt{NaN} und Bereich prüfen.
\item \textbf{Zählung ab 0}: \texttt{[1]} ist das zweite Element, \texttt{item(2)} das dritte.
\item \textbf{Whitespace wird zu Textknoten}; Attribute sind keine Kinder; Elemente haben \texttt{nodeValue} \texttt{null}.
\item \textbf{Attribut = wie im HTML, Property = wie der Browser es versteht.}
\item \textbf{\texttt{textContent} = roh, \texttt{innerText} = sichtbar}; \texttt{innerHTML} ohne, \texttt{outerHTML} mit dem Element selbst.
\item \textbf{Events blubbern nach oben}: \texttt{target} = angeklickt, \texttt{currentTarget} = wo der Listener hängt.
\item \textbf{\texttt{stopPropagation}}, wenn ein innerer Klick außen nichts auslösen soll.
\item \textbf{Im Kreis zählen mit \texttt{(i + n) \% n}}.
\item \textbf{Nur hinzufügen, nicht ändern}, wenn die Angabe es verlangt: neue Listener statt alten Code umschreiben.
\end{enumerate}
\end{merkbox}

\end{document}
